\documentclass[11pt]{amsart}
% Prefer Latin Modern; keep a Computer Modern fallback for minimal TeX setups.
\IfFileExists{lmodern.sty}{%
 \usepackage[T1]{fontenc}\usepackage{lmodern}%
}{%
 \usepackage[OT1]{fontenc}%
 \ifdefined\pdfmapfile
  \IfFileExists{pdftex.map}{}{%
   \pdfmapfile{+cm.map}%
   \pdfmapfile{+cmextra.map}%
   \pdfmapfile{+symbols.map}%
   \pdfmapfile{+euler.map}%
   \pdfmapfile{+latxfont.map}%
   \pdfmapfile{+rsfs.map}%
  }%
 \fi
}
\usepackage{amsmath,amssymb,mathtools,mathrsfs}
\usepackage{microtype}
\usepackage{enumitem}
\usepackage{needspace}
\usepackage{longtable,booktabs,array}
\usepackage{etoolbox}
\usepackage{xcolor}
\definecolor{linkblue}{RGB}{28,64,120}
\definecolor{sectionteal}{RGB}{17,103,111}
\definecolor{equationviolet}{RGB}{105,70,142}
\definecolor{citegreen}{RGB}{24,97,77}
\definecolor{urlbrown}{HTML}{A33A00}
\usepackage[colorlinks=true,linkcolor=linkblue,citecolor=citegreen,
 urlcolor=urlbrown,linktoc=all,bookmarksnumbered=true]{hyperref}
\hypersetup{pdftitle={Continuous CLD-positive potentials and mixed currents},
 pdfauthor={},pdfsubject={Local regularization and continuous mixed currents; revised version 2.6}}
\newtheorem{theorem}{Theorem}[section]
\newtheorem{proposition}[theorem]{Proposition}
\newtheorem{lemma}[theorem]{Lemma}
\newtheorem{corollary}[theorem]{Corollary}
\theoremstyle{definition}
\newtheorem{definition}[theorem]{Definition}
\theoremstyle{remark}
\newtheorem*{contactintro}{Contact rigidity}
\numberwithin{equation}{section}
\setlist[enumerate,1]{label=\textup{(\roman*)},leftmargin=2.4em}
\newcommand{\ddc}{d'd''}
\newcommand{\an}{\mathrm{an}}
\newcommand{\Ploc}{\mathcal P_{\mathrm{loc}}}
\newcommand{\Dloc}{\mathcal D_{\mathrm{loc}}}
\newcommand{\Ac}{\mathcal A_c}
\newcommand{\R}{\mathbf R}
\newcommand{\Q}{\mathbf Q}
\newcommand{\Z}{\mathbf Z}
\newcommand{\PP}{\mathbf P}
\newcommand{\Aff}{\mathbf A}
\newcommand{\OO}{\mathcal O}
\newcommand{\cX}{\mathscr X}
\newcommand{\cW}{\mathscr W}
\newcommand{\cL}{\mathscr L}
\newcommand{\cA}{\mathscr A}
\newcommand{\cE}{\mathscr E}
\newcommand{\fc}{\mathfrak c}
\newcommand{\fb}{\mathfrak b}
\newcommand{\norm}[1]{\left\lVert#1\right\rVert}
% Named results, sections, displayed equations, citations, and URLs have
% distinct, restrained link colors; the whole named reference is clickable.
\newcommand{\namedref}[2]{%
 \begingroup
 \ifstrequal{#1}{Section}{\hypersetup{linkcolor=sectionteal}}{%
  \ifstrequal{#1}{Appendix}{\hypersetup{linkcolor=sectionteal}}{}}%
 \hyperref[#2]{#1~\ref*{#2}}%
 \endgroup}
\renewcommand{\eqref}[1]{%
 \begingroup\hypersetup{linkcolor=equationviolet}%
 \hyperref[#1]{\textup{(\ref*{#1})}}%
 \endgroup}
\newcommand{\TV}{\mathrm{TV}}
\DeclareMathOperator{\supp}{supp}
\DeclareMathOperator{\Int}{Int}
\DeclareMathOperator{\divisor}{div}
\DeclareMathOperator{\im}{im}
\hyphenation{smooth}
\title[Continuous CLD-positive potentials]{Continuous CLD-positive potentials\\and mixed currents}
\author{}
\date{15 November 2026}

\begin{document}
\begin{abstract}
Let \(K\) be a non-Archimedean completion of a number field and let
\(X/K\) be smooth and quasi-projective. A finite continuous function on an
ordinary open of \(X^\an\) has positive Chambert--Loir--Ducros curvature
current if and only if it admits smooth positive approximations on a
fixed neighbourhood of each point. The proof combines restriction to
curves, deterministic complete-intersection approximation, contact
rigidity, and continuity of semipositive envelopes. Contact rigidity
identifies a continuous curve-positive metric from its contact with a
uniform Fubini--Study limit.
On a fixed regular model, sums of divisor test ideals and global generation
give one-sided \(O(m^{-1})\) Fubini--Study approximations for model-obstacle
envelopes. A single projective extension then completes local regularization. As consequences, continuous current-positive potentials admit
symmetric mixed products with any positive closed CLD current, and these
products depend only on their curvature currents. We also obtain a
projective curve criterion under a regular-model hypothesis and positivity
of the prescribed Yuan--Zhang local mixed measures.
\end{abstract}
\maketitle
\setcounter{tocdepth}{1}
\tableofcontents
\clearpage

\section{Introduction}\label{sec:introduction}

Positivity of a Chambert--Loir--Ducros current is a condition on smooth
test forms. To construct positive mixed products, one must also control
\(\ddc(uT)\) for each intermediate positive closed current \(T\).
Multiplication by a continuous potential is available, but positivity
after differentiation requires an additional argument. Local
approximation by smooth functions with positive curvature provides
that argument. The first question addressed here is whether the
current inequality itself supplies such an approximation.

\subsection{Local regularization}
Fix a number field \(k\), a non-Archimedean place \(v\), and
\(K=k_v\), with its given absolute value. Its valuation ring is \(R\).
Choose a uniformizer \(\pi\) and set \(\lambda=-\log|\pi|>0\).
All smooth forms, positive cones, and currents in this paper are those
of the first version of Chambert--Loir--Ducros \cite{CLD12}.
In particular, positivity of currents is dual to the positive cone
of \cite[Definition 5.1.1(c)]{CLD12}. An \emph{ordinary open} is open
in the Berkovich topology.

Let \(X/K\) be smooth of pure dimension \(n\), let \(O\subset X^\an\)
be an ordinary open, and let \(u:O\to\R\) be finite and continuous.
For \(n\geq1\), the condition \(\ddc[u]\geq0\) means
\begin{equation}\label{eq:current-positive}
 \langle\ddc[u],\eta\rangle
 =\int_O u\,\ddc\eta\geq0
 \quad\text{for every positive }\eta\in\Ac^{n-1,n-1}(O).
\end{equation}
Every compactly supported smooth test is allowed. For components of
different dimensions the condition is imposed componentwise, and there
is no curvature condition in dimension zero.

We write \(u\in\Ploc(O)\) if every \(x\in O\) has one ordinary
neighbourhood \(V\subset O\) on which there are smooth functions \(u_j\)
with \(\ddc u_j\geq0\) as forms and \(u_j\to u\) uniformly on compact
subsets of \(V\). The neighbourhood is fixed for the entire sequence.

\begin{theorem}[Local regularization]\label{thm:main}
Let \(K=k_v\) be a non-Archimedean completion of a number field, with
its given absolute value. Let \(X/K\) be smooth and quasi-projective,
\(O\subset X^\an\) an ordinary open, and \(u:O\to\R\) finite continuous.
Then the following are equivalent:
\begin{enumerate}
\item The original CLD curvature current \(\ddc[u]\) is positive on \(O\)
in the sense of \eqref{eq:current-positive}.
\item The function \(u\) belongs to \(\Ploc(O)\).
\end{enumerate}
After shrinking the fixed neighbourhood of each point once, the
approximating sequence can be chosen decreasing and uniformly convergent
on that neighbourhood.
\end{theorem}

The implication from smooth positive approximation to current positivity
follows by testing locally uniform convergence against \(\ddc\eta\).
The reverse implication does not assume that \(u\) is a difference of
approachable functions. The field and continuity hypotheses remain in
force throughout the theorem; in particular, the assertion does not
concern arbitrary upper-semicontinuous potentials. The proof is completed
in \namedref{Section}{subsec:regularization-proof}.

\subsection{Contact rigidity}
The projective argument uses a criterion that identifies a metric from
its contact with a semipositive lower metric. Recall that a \emph{pure
Fubini--Study metric} on an ample line bundle \(L\) is given by weights
\(m^{-1}\max_i\log|s_i|\), where the sections \(s_i\) generate \(mL\).
A \emph{uniform FS limit} is a uniform limit of such metrics.
A continuous metric is \emph{curve-positive} if its pullback by every
algebraic morphism from a smooth \(K\)-curve has positive Chern current.
These conventions are developed in \namedref{Section}{sec:metrics-curves}.

\begin{contactintro}[\namedref{Theorem}{thm:contact}]
Let \(Y/K\) be smooth, projective, and integral of dimension \(n\ge1\),
and let \(L\) be ample. Let \(p\) be a uniform FS limit on \(L\), and let
\(r\) be a finite continuous curve-positive metric on \(L\). Then
\begin{equation}\label{eq:contact-intro}
 r\ge p,\qquad \int_{Y^\an}(r-p)c_1(L,p)^n=0
 \quad\Longrightarrow\quad r=p.
\end{equation}
\end{contactintro}

The metric \(p\) need not be constructed as an envelope of \(r\).
The criterion uses the Monge--Amp\`ere measure of the lower metric and
positivity of the upper metric only on curves. It therefore applies
before a higher-dimensional positive product for \(r\) is available.
The proof in \namedref{Section}{subsec:contact-proof} passes curve Green
identities through complete-intersection approximation and then
transfers the contact identity to smooth divisors.

\subsection{Continuous mixed currents}
A current is called \emph{closed} below if both \(d'T=0\) and \(d''T=0\).
Positive currents admit multiplication by continuous functions
(\namedref{Lemma}{lem:continuous-coefficients}). Thus, for a positive closed
current \(S\), the recursive rule
\begin{equation}\label{eq:recursive-product}
 \mathcal B_S()=S,\qquad
 \mathcal B_S(u_1,\ldots,u_m)
 =\ddc\bigl(u_1\mathcal B_S(u_2,\ldots,u_m)\bigr)
\end{equation}
is meaningful whenever its intermediate currents are positive.
This is the usual Bedford--Taylor construction.

\begin{theorem}[Continuous potentials on a positive closed background]
\label{thm:background-products}
Let \(K=k_v\), let \(X/K\) be smooth quasi-projective of pure dimension
\(n\), and let \(O\subset X^\an\) be an ordinary open. Let \(S\) be an
original CLD positive current of bidegree \((p,p)\), where \(0\leq p\leq n\), with
\(d'S=d''S=0\). Suppose \(0\leq m\leq n-p\) and
\(u_1,\ldots,u_m:O\to\R\) are finite continuous functions with
\(\ddc[u_i]\geq0\). Then \eqref{eq:recursive-product} defines a positive
closed current of bidegree \((p+m,p+m)\). It has the following properties.
\begin{enumerate}
\item It is local, symmetric in the potentials, additive and positively
homogeneous in each potential and in the background.
\item If \(S_a\to S\) on every compact smooth test, where all \(S_a\)
are positive and closed of bidegree \((p,p)\), and the finite
continuous functions \(u_{i,a}:O\to\R\) converge to \(u_i\) locally
uniformly with \(\ddc[u_{i,a}]\geq0\), then
\[
 \mathcal B_{S_a}(u_{1,a},\ldots,u_{m,a})
 \longrightarrow\mathcal B_S(u_1,\ldots,u_m)
\]
on every compact smooth test.
\item Each potential enters only through its curvature current:
if \(w\) is finite continuous and \(\ddc[w]=\ddc[u_i]\), replacing
\(u_i\) by \(w\) leaves the product unchanged.
\end{enumerate}
In top degree the product is a positive Radon measure.
\end{theorem}

The recursion is described in \cite[Section 5.6.8]{CLD12}.
For continuous psh potentials, its positivity, continuity, and symmetry
are established in
\cite[Lemma 8.5.5 and Propositions 8.7.2, 8.7.4]{CLD25}.
\namedref{Theorem}{thm:main} allows us to start from current positivity alone.
Potential independence then has a short two-sided proof. If \(\ddc[h]=0\), then both \(h\) and \(-h\) satisfy
\namedref{Theorem}{thm:main}. For every positive closed \(T\) this gives
\begin{equation}\label{eq:two-sided-intro}
 \ddc(hT)\geq0,\qquad -\ddc(hT)=\ddc((-h)T)\geq0,
 \qquad\text{hence }\ddc(hT)=0.
\end{equation}
No representation of \(h\) as a piecewise-linear or logarithmic function
is required. \namedref{Theorem}{thm:background-products} concerns continuous
potentials; it does not assert a monotone convergence theorem for
bounded discontinuous ones. The theorem is proved in
\namedref{Section}{subsec:mixed-current-proof}.

\subsection{Adelic mixed measures}
For the arithmetic application, use the rational boundary completion of
\cite[Sections 2.4--2.5]{YZ}. Its objects retain the specified
line-bundle identifications. A strongly nef object has a Cauchy
presentation by nef model objects, and an integrable object is a
difference of two strongly nef objects, with rational tensor powers
allowed. The order defining the completion compares the full
Cartier-and-metric boundary datum, as recalled in
\namedref{Section}{sec:adelic}.

\begin{theorem}[Positivity of the prescribed adelic mixed measure]
\label{thm:adelic}
Let \(k\) be a number field, \(v\) a non-Archimedean place, and \(K=k_v\).
Let \(U/k\) be smooth quasi-projective of pure dimension \(n\geq1\).
Let \(\overline L_1,\ldots,\overline L_n\) be rational integrable adelic
line bundles with their specified identifications in the boundary
completion of \cite[Sections 2.4--2.5]{YZ}. Suppose their local weights
at \(v\) are finite continuous and have positive original CLD curvature
currents. Then the signed measure prescribed in
\cite[Sections 3.6.7--3.6.8]{YZ} is positive:
\[
 c_1(\overline L_1)_v\cdots c_1(\overline L_n)_v\geq0
 \quad\text{on }U_K^\an.
\]
\end{theorem}

The signed measure is defined by expanding integrable presentations.
Its positivity follows by identifying that expansion with the local
positive product of the target weights. This conclusion concerns the
given measure and does not assert arithmetic nefness of the factors.
The comparison is given in \namedref{Section}{subsec:adelic-proof}.

\subsection{Relation to earlier work}
Chambert--Loir--Ducros distinguish distributional positivity from local
psh approachability in \cite[Definition 5.6.3 and Remark 5.6.4]{CLD12};
their stable products are treated in
\cite[Corollary 5.6.5 and Section 5.6.8]{CLD12}.
The later continuous psh class is defined by uniform closure
\cite[Definition 8.1.7]{CLD25}. We do not identify a continuous
current-positive function with a member of that class before proving
\namedref{Theorem}{thm:main}.

On curves, Wanner compares CLD and Thuillier subharmonicity
\cite[Corollaries 4.8 and 5.4]{Wanner}.
Gubler--Jell--Rabinoff provide descent, harmonic tropicalizations,
and smooth approximation \cite{GJRforms,GJRcurves}.
The finite residue field is important here: weakly smooth functions
on the relevant curves are CLD-smooth. \namedref{Section}{sec:metrics-curves}
retains the passage from the original current inequality over \(K\)
to those curve results.

Torus-invariant pluripotential theory is developed in \cite{BGJK}.
Piecewise-linear semipositivity and positivity tested on curves are
studied in \cite{GRpl,GRpsh}, and Wu compares local approximation
classes in \cite[Appendix A]{Wu}. Our starting hypothesis is instead
the original current inequality \eqref{eq:current-positive} for an
arbitrary finite continuous potential on a smooth variety.

The envelope argument follows the test-ideal method of
Fang--Gubler--K\"unnemann \cite{FGK}, using the mixed-characteristic
test ideals of \cite{BMP}. We use test ideals of section divisors directly.
The comparison on a higher normal model allows an
obstacle on a higher normal model while subadditivity is applied on
a fixed regular model. It is not a direct application of
\cite[Theorems 5.1 and 7.1]{FGK}. Orthogonality, once continuity is
established, is \cite[Theorem 6.7]{BGM}.
The final adelic application identifies a measure already constructed
by Yuan--Zhang; it does not introduce another mixed measure.

\subsection{Proof strategy}
The proof separates restriction, contact, and envelopes so that none
uses regularization in a form still to be established.

\emph{Step 1: restrict the current inequality and extend locally.}
In a normal disc, let
\(\beta_N=\ddc R_{2,\varepsilon}(\log|t|,-N)\), where
\(R_{2,\varepsilon}\) is a smooth convex regularized maximum.
For the zero section \(i:B\hookrightarrow B\times D_s\),
Poincar\'e--Lelong and continuous multiplication give
\begin{equation}\label{eq:normal-limit-intro}
 \beta_N\longrightarrow i_*[B],
 \qquad u\beta_N\longrightarrow i_*[u|_B].
\end{equation}
Testing against pullbacks from \(B\) proves positivity on the zero
section. The \'etale normal form of a smooth pair and Berkovich's
theorem on \'etale germs then give restriction to smooth
subvarieties. Curve regularization and removal of finitely many
closed points prove positivity after every algebraic curve map.
An algebraic \'etale presentation also supplies an ambient analytic
retraction, extending \(u\) to a continuous curve-positive function
\(v\) near the original point
(\namedref{Section}{sec:restriction}).

\emph{Step 2: approximate mixed currents by curves and prove contact.}
For a generating tuple \(s_0,\ldots,s_b\), choose smooth cuts of
the form
\[
 \sigma_N=\sum_{i=0}^b a_{i,N}s_i^N,
 \qquad a_{i,N}\in\pi^i(1+\pi R).
\]
On a tropical presentation of any fixed smooth test, equality of
the norms of two terms lies on an affine hyperplane, apart from
finitely many exceptional \(N\). Distinct coefficient valuations
exclude ties that persist for all \(N\).
Absolute tropical integration therefore gives an eventual
\(O(N^{-1})\) error on that test. Successive cuts and weak closure
yield normalized individual smooth complete-intersection curves,
with fixed intersection degrees
(\namedref{Proposition}{prop:curve-approximation}).

For the contact criterion \eqref{eq:contact-intro}, put
\(g=r-p\) and \(\alpha=c_1(L,p)\).
Passing curve Green identities through the approximation of
\(T=\alpha^{n-1}\) gives
\[
 \nu=\ddc(gT)=\sigma-\alpha^n,\qquad \sigma\geq0,\qquad
 0\leq-\int g\,d\nu=-\int g\,d\sigma\leq0.
\]
The nonnegative energy of \(g+th\), for smooth \(h\) and \(t\in\R\),
then forces \(\nu=0\).
For a smooth divisor \(D=\divisor(s)\) in \(mL\), set
\(\rho_N=\max\{mp,\log|s|+N\}\). Fixed-potential continuity gives
\begin{equation}\label{eq:divisor-intro}
 c_1(mL,\rho_N)\wedge\alpha^{n-1}
 \longrightarrow i_*\bigl(c_1(L|_D,p|_D)^{n-1}\bigr).
\end{equation}
Green symmetry transfers contact to \(D\), and induction on
dimension proves \eqref{eq:contact-intro}
(\namedref{Section}{sec:contact}). No mixed product with the upper metric
\(r\) is assumed in this argument.

\emph{Step 3: construct continuous envelopes on a fixed regular model.}
Represent a model obstacle on a higher normal model \(\cW\) as
\(r=\theta+\log|\cE|\), with \(\cE\) an invertible vertical ideal.
For each admissible section \(s\in H^0(Y,lL)\), let \(D_s\) be its
divisor on a fixed regular model \(\cX\). We take the sum
\[
 \mathfrak j_m=\sum_{l\ge1}\sum_{0\ne s\in V_l}
       \tau_{\cX}\!\left(\frac mlD_s\right),
 \qquad V_l=\{s:\|s\|_{lr,\sup}\le1\}.
\]
Subadditivity gives the lower bound below. Birational comparison,
followed by clearing a fixed discrepancy denominator, gives the upper bound:
\begin{equation}\label{eq:ideal-intro}
 \pi^B\mathfrak j_m\OO_{\cW}\subseteq\cE^m,\qquad
 P_L(r)\leq\theta+\frac1m\log|\mathfrak j_m|
 \leq r+\frac{B\lambda}{m}.
\end{equation}
Global generation and one fixed auxiliary line bundle turn these bounds
into pure-FS approximations of \(P_L(r)\) on \(L\).
On \(\cX=\PP_R^a\), the identity
\(K_{\cX}+(a+1)\OO(1)=0\) cancels the auxiliary twist directly and gives
\begin{equation}\label{eq:rate-intro}
 0\leq q_m-P_L(r)\leq\frac{B\lambda}{m}.
\end{equation}
Model-metric density and the one-Lipschitz property of the envelope
treat continuous obstacles. Orthogonality \cite[Theorem 6.7]{BGM}
then applies (\namedref{Section}{subsec:envelope-proof}).

\emph{Step 4: extend once to projective space and apply contact rigidity.}
After adding a constant and shrinking the ambient neighbourhood, arrange
\(0<v+c<\lambda/2\). Choose one compact polynomial norm box \(C\) and
an FS metric \(q\), with a section \(s_0\), such that
\[
 b=-\log\|s_0\|_q,\qquad b=0\text{ near }x,
 \qquad\{b\le\lambda\}=C.
\]
On \(\{b<\lambda\}\), use the weight \(\max\{v+c,b\}\); on
\(\{b>\lambda/2\}\), use \(q\). The two definitions agree on the
overlap and give a global curve-positive metric \(r\)
(\namedref{Lemma}{lem:projective-extension}). Its envelope \(p=P_L(r)\)
satisfies the contact identity, so \namedref{Theorem}{thm:contact} gives
\(r=p\). Smooth FS approximation and restriction to one fixed
neighbourhood recover \(u\). Explicit constant translations make the
approximants decreasing (\namedref{Section}{subsec:regularization-proof}).
The same envelope--contact argument gives the projective curve criterion
under the retained regular-model hypothesis
(\namedref{Corollary}{cor:projective-criterion}).

\subsection{Organization}
\namedref{Section}{sec:coefficients} proves the continuous-coefficient
estimate and the continuity of a fixed approachable potential acting
on positive closed currents. \namedref{Section}{sec:metrics-curves} supplies
smooth Fubini--Study approximation, the measure normalization, and
the Green identities on curves. \namedref{Section}{sec:restriction}
establishes restriction and the ambient extension.
\namedref{Section}{sec:contact} combines complete-intersection approximation
with the direct divisor limit to prove contact rigidity.
\namedref{Section}{sec:envelopes} proves envelope continuity on a fixed regular
model, with projective space as a corollary. \namedref{Section}{sec:local-proof}
constructs the projective extension and completes local regularization. The mixed-current and curve-criterion consequences are
proved in \namedref{Section}{sec:consequences}; \namedref{Section}{sec:adelic}
identifies the prescribed adelic measure.
\namedref{Appendix}{app:comparison} records a quantitative comparison of
mixed measures, independently of the contact argument.

\section{Continuous coefficients and approachable potentials}
\label{sec:coefficients}

We use the smooth cutoff and density results
\cite[Corollary 3.3.4 and Propositions 3.3.5--3.3.6]{CLD12},
the Poincar\'e--Lelong formula \cite[Theorem 4.6.5]{CLD12}, and
uniform stability of products
\cite[Corollary 5.6.5]{CLD12}. Weak convergence of currents always
means convergence on every compactly supported smooth test.
Convergence of measures on continuous tests will be stated explicitly.

\begin{definition}\label{def:approachable}
On an ordinary open \(O\subset X^\an\), a smooth function is locally
of the form
\[
 F(\log|f_1|,\ldots,\log|f_r|),
\]
where the \(f_i\) are analytic units and \(F\) is smooth near their
logarithmic image. The class \(\Ploc(O)\) is defined in
\namedref{Section}{sec:introduction}. We write \(u\in\Dloc(O)\) if \(u\)
is locally a difference of two members of \(\Ploc\).
For metrics, these conditions apply to their weights in local frames.
The weight in a frame \(e\) is \(-\log\|e\|\).
\end{definition}

A change of frame changes a weight by the logarithm of the norm of an
analytic unit. That function is smooth and has zero \(d'd''\).
Continuous multiplication of positive currents is known from
\cite[Corollary 5.4.7]{CLD12} and \cite[Lemma 8.5.1]{CLD25}.
We give the coefficient construction because it also controls
weakly convergent nets.

\Needspace{14\baselineskip}
\begin{lemma}[Continuous coefficients]\label{lem:continuous-coefficients}
Let \(O\) be an ordinary open in a smooth pure \(n\)-dimensional
\(K\)-variety. Let \(T_a\) be positive currents of bidegree
\((n-q,n-q)\), \(0\leq q\leq n\), converging weakly to \(T\).
Then \(T\) is positive. For \(S=T\) and for each \(S=T_a\), every finite
continuous \(f:O\to\R\) determines a unique current \(fS\) by
\[
 \langle fS,\eta\rangle
   =\lim_{j\to\infty}\langle S,f_j\eta\rangle,
\]
where \(\eta\) is any compactly supported smooth test and \(f_j\)
are smooth functions converging uniformly to \(f\) near
\(\supp\eta\). The limit exists and is independent of the choice
of \(f_j\).

On a fixed compact test chart, the scalar coefficient measures of
\(T_a\) have eventually uniformly bounded total variations. The same
holds after wedging \(T_a\) with any fixed smooth form, in the
bidegrees where the wedge is defined. The cone of positive currents
is pointed. If \(f_a\to f\) locally uniformly, with \(f_a,f\)
finite continuous, then
\begin{equation}\label{eq:moving-coefficients}
 f_aT_a\longrightarrow fT
\end{equation}
on every compact smooth test. Nets are allowed.
\end{lemma}

\begin{proof}
Positivity of \(T\) follows by taking the limit on each positive test.
We establish the coefficient estimate locally and then use a finite
smooth partition. Fix a moment chart with \(r\) moment coordinates,
a compactly supported \((q,q)\)-form \(\eta\), and a nonnegative
smooth cutoff \(\chi\) equal to one near its support. Let \(J\) denote
the involution interchanging the two differential factors, and write
\[
 \eta_{\mathrm s}=\tfrac12\bigl(\eta+(-1)^qJ\eta\bigr),
 \qquad \eta_{\mathrm a}=\eta-\eta_{\mathrm s}.
\]
Wedging \(\eta_{\mathrm a}\) with \(\gamma\wedge J\gamma\), for an
\((r-q,0)\)-form \(\gamma\), gives zero. In the convention
\cite[Definition 5.1.1(c)]{CLD12}, both signs of
\(\eta_{\mathrm a}\) are therefore positive. A positive current
annihilates \(\eta_{\mathrm a}\), also after multiplication by a
smooth real scalar.

For symmetric forms the positivity test is a finite-dimensional
positive-semidefinite matrix condition. The form
\[
 \Omega=\frac1{q!}\left(\sum_{j=1}^r d'x_j\wedge d''x_j\right)^q
 \quad(\Omega=1\text{ if }q=0)
\]
has a positive-definite coefficient matrix. Boundedness of the
coefficient matrix of \(\eta_{\mathrm s}\) on the fixed compact set
gives \(C_\eta\) such that, for every smooth real \(h\),
\[
 C_\eta\|h\|_{\supp\eta}\chi\Omega\pm h\eta_{\mathrm s}\geq0.
\]
If \(h\) uses additional moment coordinates, pull this inequality
to a common refinement; positivity is preserved. The forms extend
by zero outside the chart. For \(S=T_a\) or \(T\), we obtain
\begin{equation}\label{eq:coefficient-estimate}
 |\langle S,h\eta\rangle|
 \leq C_\eta\|h\|_{\supp\eta}\langle S,\chi\Omega\rangle.
\end{equation}
At this stage \(h\) is smooth, so the pairing is already defined.

The right-hand evaluations for \(T_a\) converge and are eventually
bounded. Apply \eqref{eq:coefficient-estimate} to the finitely many
coordinate test forms. Their scalar functionals extend uniquely to
continuous compactly supported functions, and the Riesz theorem
represents them by signed Radon measures. The same estimate bounds
their total variations. These extensions agree on overlaps by
uniform smooth density. On a fixed test stage \(C_\eta\) is bounded
by a constant times the supremum norm of the coefficient matrix,
so the extended functionals define currents.

We can now define \(fS\), for \(S=T\) or \(S=T_a\), by multiplying
these coefficient measures by \(f\). Smooth density supplies
the functions \(f_j\) in the statement; estimate
\eqref{eq:coefficient-estimate} proves that their pairings converge
to this product independently of the approximation. The same local
estimate proves continuity on each smooth test stage.

For fixed continuous \(f\), choose smooth \(f_j\to f\) uniformly
on an enlarged compact set. For an eventually uniform \(M_\eta\),
\[
 |\langle fT_a-fT,\eta\rangle|
 \leq|\langle T_a-T,f_j\eta\rangle|
       +M_\eta\|f-f_j\|_\infty.
\]
First fix \(j\), let \(a\) tend to its limit, and then let \(j\)
increase. This proves \(fT_a\to fT\).
For varying \(f_a\), the additional term satisfies
\[
 |\langle(f_a-f)T_a,\eta\rangle|
 \leq M_\eta\|f_a-f\|_{\supp\eta}\longrightarrow0.
\]
A finite partition treats any compact test. The coefficient bounds
also control wedges with a fixed smooth form, whose coefficients are
bounded on compact charts. Finally, if both \(S\) and \(-S\) are
positive, their evaluations on \(\chi\Omega\) vanish; applying
\eqref{eq:coefficient-estimate} to \(S\) gives \(S=0\).
\end{proof}

\begin{proposition}[A fixed approachable potential]\label{prop:fixed-potential}
Let \(T\) be a positive current on \(O\) with \(d'T=d''T=0\), and let
\(u\in\Ploc(O)\). Whenever the bidegrees permit it,
\[
 M_u(T):=\ddc(uT)
\]
is positive and closed. For fixed \(u\), the operation \(M_u\) is
continuous on the cone of positive closed currents for weak
convergence. The same is true jointly when \(u_a\to u\) locally
uniformly and all \(u_a,u\) belong to \(\Ploc\).
Successive such operations are symmetric in their potentials.
\end{proposition}

\begin{proof}
On a fixed smaller neighbourhood choose smooth \(u_j\to u\) with
\(\ddc u_j\geq0\). Closedness gives
\[
 \ddc(u_jT)=\ddc u_j\wedge T\geq0.
\]
The curvature \(\ddc u_j\) is a symmetric positive \((1,1)\)-form,
hence strongly positive, so its wedge preserves positivity. \namedref{Lemma}{lem:continuous-coefficients}
and continuity of differentiation give the positive limit
\(\ddc(uT)\). Both differentials annihilate this limit.
Equation~\eqref{eq:moving-coefficients} proves the two continuity
assertions. Induction gives joint continuity for successive
operations. Simultaneous smooth positive approximations of finitely
many potentials then give symmetry from the smooth wedge product.
\end{proof}

For a continuously metrized line bundle with approachable weights, write
\(c_1(L,p)\wedge T\) for this operation in local frames. The local
definitions agree because a transition logarithm \(\ell\) is smooth,
\(\ddc\ell=0\), and \(\ddc(\ell T)=0\).

\section{Metrics and Green identities on curves}\label{sec:metrics-curves}

\subsection{Fubini--Study metrics and normalization}
\label{app:maxima}% Historical anchor for the relocated standard construction.
\begin{definition}\label{def:fs}
Let \(L\) be an ample rational line bundle on a projective \(Y/K\).
A \emph{pure Fubini--Study metric} on \(L\) is defined by an integer
\(d>0\), with \(dL\) an actual line bundle, and a generating tuple
\(s_0,\ldots,s_b\in H^0(Y,dL)\). Its weight in a local frame is
\[
 q_e=\frac1d\max_{0\leq j\leq b}\log|s_j/e^{\otimes d}|.
\]
Coefficients multiplying the sections belong to \(K\). A
\emph{uniform FS limit} is a uniform limit of these metrics.
A metric is \emph{globally smoothly approximable} if it is a uniform
limit of global smooth metrics with positive curvature forms.
\end{definition}

A model metric is \emph{nef} if, after clearing tensor denominators,
its model line bundle has nonnegative degree on every integral curve
of the special fibre. A continuous metric is \emph{semipositive}
if it is a uniform limit of nef model metrics. For an ample rational
line, this is equivalent to being a uniform pure-FS limit
\cite[Corollary 7.6 and Theorem 7.8]{BE}. All uniform differences of
metrics are scalar differences on the same specified line bundle.

We use the standard regularized maximum from
\cite[proof of Proposition 6.3.2]{CLD12}.

\begin{proposition}[Regularized maximum]\label{prop:regularized-max}
Fix an even nonnegative \(C^\infty\) function \(\rho\), supported
in \([-1,1]\), with integral one. For \(h\geq1\) and
\(\varepsilon>0\), put
\[
 R_{h,\varepsilon}(t)
 =\int_{\R^h}\max_i(t_i+\varepsilon z_i)
                    \prod_{i=1}^h\rho(z_i)\,dz.
\]
This function is smooth, convex, and nondecreasing in each variable.
It satisfies
\[
 \max_i t_i\leq R_{h,\varepsilon}(t)\leq\max_i t_i+\varepsilon,
 \qquad
 R_{h,\varepsilon}(t+c\mathbf1)=R_{h,\varepsilon}(t)+c.
\]
An entry more than \(2\varepsilon\) below the maximum is inactive
on a neighbourhood of that point. If smooth \(a_1,\ldots,a_h\)
have positive curvature, so does \(R_{h,\varepsilon}(a_1,\ldots,a_h)\).
\end{proposition}

Translation covariance gives the metric transition law, and the
inactive-coordinate property handles zeros of the generating sections.
These are precisely the features needed in the next lemma.

\begin{lemma}[Smooth approximation of FS metrics]\label{lem:smooth-fs}
A pure-FS metric, and hence a uniform FS limit, is globally smoothly
approximable. If \(p\) is globally smoothly approximable on \(L\),
if \(m>0\) is an integer for which \(mL\) is a line bundle, and if
\(s\in H^0(Y,mL)\), then for each \(N\in\R\) the metric
\[
 \rho_N=\max\{mp,\log|s|+N\}
\]
on \(mL\) is globally smoothly approximable.
\end{lemma}

\begin{proof}
For a pure-FS metric from sections of \(dL\), pull back the smooth
positive approximations of the canonical metric on
\(\OO_{\PP^b}(1)\) from \cite[Proposition 6.3.2]{CLD12}, and
divide their weights by \(d\). A diagonal choice treats uniform
FS limits.

The second assertion is immediate for \(s=0\). Otherwise, take
smooth positive \(p_j\to p\) and regularize \(\max\{mp_j,\log|s|+N\}\).
The logarithmic entry is inactive near the divisor of \(s\);
elsewhere it is smooth with zero curvature. The resulting global
metrics are positive and have error at most
\(m\|p_j-p\|_\infty+\varepsilon_j\).
This is the same regularized-maximum construction, applied to a
smooth weight and a section weight.
\end{proof}

\begin{lemma}[Degree and comparison of measures]\label{lem:degree}
On a smooth projective \(Y/K\) of dimension \(n\), globally smoothly
approximable metrics on \(L_1,\ldots,L_n\) have total mixed mass
\(\deg_K(L_1\cdots L_n)\). For uniform FS limits, the model-theoretic
mixed measure agrees with the CLD product.
\end{lemma}

The mass assertion is \cite[Corollary 6.4.4]{CLD12}.
The model comparison is
\cite[Theorem 6.9.3 and Proposition 6.9.7]{CLD12}, with the
discretely valued normalization of \cite[Section 6.9.6]{CLD12}.
For uniform FS limits, use uniform continuity of the model
measures \cite[Proposition 2.23(d) and Remark 2.24]{BGM} and of
the approachable CLD products. Rational tensor denominators are
cleared and divided on both sides.

Degrees in this paper are over the stated ground field. A variety
over a finite extension \(F/K\), regarded as a \(K\)-variety,
contributes the structural factor \([F:K]\) to intersection and
integration. Indeed, a moment map followed by
\((\mathbf G_m^r)_F\to(\mathbf G_m^r)_K\) has that degree.
This operation is distinct from scalar extension of a \(K\)-variety:
the normalized scalar-extension integration formula used next has
no extra degree factor.

\begin{definition}\label{def:curve-positive}
A continuous function on an ordinary algebraic analytic open is
\emph{curve-positive} if its pullback by every algebraic morphism
from a smooth \(K\)-curve has positive CLD curvature current on the
inverse-image open. For a continuous metric the same condition is
imposed on the pulled-back line bundle. Curves may be disconnected.
The definition uses algebraic curve maps over the given \(K\);
it does not build positivity after arbitrary field extension into
the hypothesis.
\end{definition}

\subsection{Regularization and energy in dimension one}
\begin{proposition}[Curve regularization]\label{prop:curves}
Let \(C/K\) be a smooth algebraic curve and \(W\subset C^\an\)
an ordinary open. A finite continuous CLD-positive function on \(W\)
belongs to \(\Ploc(W)\). A finite continuous function with zero
Laplacian on a connected proper smooth curve is constant.
\end{proposition}

\begin{proof}
For a finite Galois extension \(F/K\), average a nonnegative compact
smooth scalar test \(\eta\) on \(W_F\) over the Galois group.
The invariant average descends by
\cite[Proposition 14.2]{GJRforms}. It is smooth because weakly
smooth and CLD-smooth forms coincide on these curves
\cite[Remark 15.10]{GJRforms}: their residue fields are algebraic
over a finite field. If \(\overline\eta\) is the descended test,
the scalar-extension integration formula
\cite[Proposition 14.5]{GJRforms} gives
\[
 \int_{W_F}u_F\,\ddc\eta=\int_Wu\,\ddc\overline\eta.
\]
Uniform smooth density permits the continuous coefficient \(u\).
Consequently current positivity passes to finite Galois extensions.

Let \(Y\) be a compact strictly affinoid curve domain inside \(W\).
After such an extension choose a semistable skeleton \(\Gamma\)
with retraction \(\tau\), edge lengths, and positive weights \(w_e\).
For an edgewise smooth \(\varphi\), constant near vertices and zero
near boundary vertices, \(\varphi\circ\tau\), extended by zero,
is a compact smooth test. Its Laplacian has no vertex atoms, and
integration on the annular edges gives
\begin{equation}\label{eq:curve-skeleton}
 \int_\Gamma\varphi\,d\tau_*(\ddc[u]|_Y)
 =\sum_e w_e\int_e u|_e\,\varphi_e''\,ds.
\end{equation}
This is also the retraction comparison
\cite[Proposition 6.3.2]{GJRcurves}; it does not assume that \(u\)
factors through \(\tau\).

Tests inside an edge show that \(u|_e\) is convex. The positive
Radon measure \(\ddc[u]\) has finite mass on the compact domain.
Thus its edge second derivatives have finite mass and finite
one-sided derivatives at the endpoints. Integration by parts in
\eqref{eq:curve-skeleton}, allowing arbitrary constants near the
interior vertices, gives a nonnegative weighted sum of outgoing
slopes at each such vertex. The finite-graph maximum principle
therefore applies to \(u\).

To pass from graph inequalities to subharmonicity, we follow the
harmonic-comparison argument of \cite[Theorem 4.7]{Wanner}.
A harmonic comparison function becomes skeletal after semistable
refinement \cite[Section 15.6 and Proposition 15.7]{GJRforms}.
Its boundary comparison with \(u\) holds on that skeleton.
Any prescribed finite set of type-II points can be included in a
refinement after a further finite Galois extension
\cite[Remark 5.4.2]{GJRcurves}. The boundary and harmonic
compatibilities are \cite[Lemma 5.4.1 and Proposition 5.6.8]{GJRcurves}.
Density of these points and continuity prove harmonic comparison
on all of \(Y\). Hence \(u\) is Thuillier-subharmonic.

By \cite[Proposition 15.11]{GJRforms}, \(u\) has a decreasing net
of weakly smooth psh approximants on a fixed relatively compact
open. They are CLD-smooth by the finite-residue-field comparison.
On a smaller compact neighbourhood, Dini's finite-cover argument
selects a uniformly convergent sequence. On a curve, positivity of
the curvature current of a smooth function implies positivity of
its curvature form: a negative coefficient on an open
one-dimensional tropical cell is detected by a nonnegative scalar
test there; continuity handles faces. This proves approachability.

For a zero-Laplacian function on a proper curve,
\eqref{eq:curve-skeleton} makes its restriction to every refined
skeleton harmonic. A harmonic function on a connected compact
graph without boundary is constant. Refinement through arbitrary
type-II points and descent give the final assertion.
\end{proof}

The smooth Green identity is \cite[Theorem 3.12.2]{CLD12}.
We extend it and the smooth chain rule to the continuous gaps
below using curve regularization and bounded Laplacian variation.

\begin{lemma}[Green identities on curves]\label{lem:curve-green}
Let \(r,p\) be finite continuous metrics with positive Chern currents
on the same rational line \(L\) on a smooth proper curve \(C/K\).
Their gap \(g=r-p\) has smooth approximants \(a_j\to g\) uniformly
with \(\sup_j\|\ddc a_j\|_{\TV}<\infty\). The two Chern measures
have mass \(\deg_KL\), so \(\|\ddc g\|_{\TV}\leq2\deg_KL\).

On every finite real span \(E\) of such gaps and smooth functions,
\begin{equation}\label{eq:curve-green}
 \int_C a\,\ddc b=\int_C b\,\ddc a,\qquad
 Q_C(a):=-\int_C a\,\ddc a\geq0 \quad(a,b\in E).
\end{equation}
Moreover,
\[
 \Gamma_C(a)=\ddc(a^2/2)-a\,\ddc a
\]
is a positive measure of mass \(Q_C(a)\), and
\[
 \ddc\Phi(a)=\Phi'(a)\ddc a+\Phi''(a)\Gamma_C(a)
\]
for every smooth real \(\Phi\) defined near the range of \(a\).
If \(Q_C(a)=0\), then \(\ddc a=0\) and \(a\) is constant on each
connected component.
\end{lemma}

\begin{proof}
In finitely many fixed frame neighbourhoods, use
\namedref{Proposition}{prop:curves} to approximate the weights by smooth
subharmonic \(r_{i,j},p_{i,j}\). For a compact \(F\) in one such
neighbourhood, choose a nonnegative smooth cutoff \(\xi=1\) near \(F\).
Then
\[
 (\ddc r_{i,j})(F)
 \leq\int\xi\,\ddc r_{i,j}
 =\int r_{i,j}\,\ddc\xi
 \leq\|r_{i,j}\|_{\supp\ddc\xi}\|\ddc\xi\|_{\TV}.
\]
This bounds the local Laplacians of
\(g_{i,j}=r_{i,j}-p_{i,j}\). For a compactly supported
\(\rho\geq0\), the square identity gives
\[
 \int\rho\,d'g_{i,j}\wedge d''g_{i,j}
 =\int(g_{i,j}^2/2)\ddc\rho
   -\int\rho g_{i,j}\ddc g_{i,j},
\]
so the local energies are bounded as well.

Take a finite smooth partition \(\chi_i\) and set
\(a_j=\sum_i\chi_i g_{i,j}\). Cauchy--Schwarz for integration on
the curve bounds the cross terms by
\[
 \|d'\chi_i\wedge d''g_{i,j}
       +d'g_{i,j}\wedge d''\chi_i\|_{\TV}
 \leq2Q_C(\chi_i)^{1/2}
 \left(\int\rho_i\,d'g_{i,j}\wedge d''g_{i,j}\right)^{1/2},
\]
where \(\rho_i=1\) near \(\supp\chi_i\).
The product rule therefore gives the uniform global variation bound.

The mass formula is \cite[Corollary 18.1.7]{CLD25}.
Uniform convergence and the variation bound imply weak convergence
of \(\ddc a_j\) on continuous tests and
\[
 -\int a_j\ddc a_j\longrightarrow-\int a\ddc a.
\]
The cited smooth Green formula and
\(\ddc(a_j^2/2)-a_j\ddc a_j=d'a_j\wedge d''a_j\) pass to the limit.
The latter positive measures have convergent masses. The smooth
chain rule passes as well. Choosing approximations simultaneously
proves these assertions on a fixed finite span \(E\).
If \(Q_C(a)=0\), the positive measure \(\Gamma_C(a)\) is zero.
The nonnegative quadratic \(Q_C(a+th)\), for smooth \(h\) and
arbitrary \(t\in\R\), has zero linear coefficient. By Green symmetry,
\(\int h\,\ddc a=0\) for every such \(h\). Thus \(\ddc a=0\), and
\namedref{Proposition}{prop:curves} gives constancy.
\end{proof}

\section{Restriction and an ambient extension}\label{sec:restriction}

Restriction of a current to a subvariety is not a formal operation.
For continuous potentials it follows from a positive kernel in the
normal direction.

\begin{lemma}[A normal kernel]\label{lem:normal-kernel}
Let \(B\) be an ordinary open in the analytification of a smooth
\(K\)-variety of pure dimension \(d\geq1\), let
\(D_s=\{|t|<s\}\), and let \(i:B\hookrightarrow B\times_KD_s\)
be the zero section. For a fixed \(\varepsilon>0\), put
\[
 v_N=R_{2,\varepsilon}(\log|t|,-N),\qquad \beta_N=\ddc v_N.
\]
Then, for every finite continuous \(u\) on the product,
\[
 \beta_N\longrightarrow i_*[B],\qquad
 u\beta_N\longrightarrow i_*[u|_B]
\]
as currents. In particular, \(\ddc[u]\geq0\) implies
\(\ddc[u|_B]\geq0\).
\end{lemma}

\begin{proof}
The form \(\beta_N\) is smooth and strongly positive, with support in
\[
 e^{-N-2\varepsilon}\leq|t|\leq e^{-N+2\varepsilon}.
\]
The logarithm is locally integrable. Dominated convergence and
Poincar\'e--Lelong give
\(\ddc[v_N]\to\ddc[\log|t|]=i_*[B]\), with multiplicity one.
\namedref{Lemma}{lem:continuous-coefficients} permits the fixed multiplier \(u\).
For a positive \(\eta\in\Ac^{d-1,d-1}(B)\),
\(\operatorname{pr}^*\eta\wedge\beta_N\) is a positive compact test
for large \(N\). Indeed, its support lies above a fixed compact
subset of \(B\) and inside a fixed closed normal disc. Thus
\[
 \begin{aligned}
 0&\leq\langle\ddc[u],\operatorname{pr}^*\eta\wedge\beta_N\rangle\\
  &=\langle u\beta_N,\operatorname{pr}^*(\ddc\eta)\rangle
    \longrightarrow\langle\ddc[u|_B],\eta\rangle.
 \end{aligned}
\]
Fixed base and normal cutoffs make the middle expression a compact
test pairing without altering it.
\end{proof}

\begin{lemma}[A local product neighbourhood]\label{lem:tubular}
Let \(i:Z\hookrightarrow Y\) be a closed immersion of smooth
\(K\)-varieties, let \(x\in Z^\an\), and let \(r\) be the codimension
near \(x\). There are ordinary neighbourhoods \(V\subset Z^\an\)
and \(W\subset Y^\an\), a radius \(s>0\), and an analytic isomorphism
\[
 \Phi:V\times_KD_s^r\xrightarrow{\sim}W
\]
whose zero section is \(i|_V\). Consequently restriction of a finite
continuous positive potential along \(i\) preserves current positivity.
\end{lemma}

\begin{proof}
Use the \'etale local form of a smooth pair
\cite[Lemma 37.17.1, Tag 0FUE]{Stacks}. After Zariski shrinking there
is an \'etale map
\[
 \Psi=(z,F):Y\longrightarrow\Aff_K^{d+r},\qquad
 Z=\Psi^{-1}(\Aff_K^d\times\{0\}),
\]
where \(d=\dim Z\) and \(z|_Z\) is \'etale. Form
\(E=Y\times_{\Aff_K^d}Z\). The two maps
\[
 E\xrightarrow{\,p\,}Y,\qquad
 E\xrightarrow{\,q=(\operatorname{pr}_Z,F\circ p)\,}
                       Z\times\Aff_K^r
\]
are \'etale by base change. At the diagonal point \(e\) above \(x\),
both induce the identity on the completed residue field
\(\mathscr H(x)\). Berkovich's equivalence for \'etale germs
\cite[Theorem 3.4.1]{Ber93} therefore makes both analytic germs
isomorphisms. By the definition of these germs
\cite[Section 3.4]{Ber93}, they are represented by isomorphisms of
ordinary open neighbourhoods. Their composition identifies a
neighbourhood of \(x\) in \(Y^\an\) with an ordinary open
\(G\subset Z^\an\times\Aff_K^{r,\an}\), matching \(Z\) with the
zero section.

Choose a compact affinoid neighbourhood \(V_0\) of \(x\) in \(Z^\an\)
whose zero section lies in \(G\). In the product of \(V_0\) with a
fixed closed normal polydisc, the complement of \(G\) is compact
and disjoint from the zero section. If nonempty, it has a positive
minimum of \(\max_i|t_i|\). Choose \(s\) smaller than the initial
polydisc radius and, when the complement is nonempty, smaller than
that minimum. Then \(V\times D_s^r\subset G\), where \(V=\Int V_0\).
Its image is the required ordinary open \(W\).
Finally apply \namedref{Lemma}{lem:normal-kernel} in each normal coordinate;
a zero-dimensional target has no curvature condition.
\end{proof}

\begin{proposition}[Curve maps and extension]\label{prop:curve-restriction}
Let \(u\) be a finite continuous CLD-positive function on an ordinary
open \(O\subset X^\an\), where \(X/K\) is smooth. Then:
\begin{enumerate}
\item For every algebraic morphism \(f:C\to X\) from a smooth
\(K\)-curve, \(u\circ f^\an\) is CLD-positive on
\((f^\an)^{-1}(O)\).
\item Near each point of \(O\), an affine embedding of \(X\) has
a finite continuous curve-positive ambient extension \(v\) of \(u\).
\end{enumerate}
\end{proposition}

\begin{proof}
Argue componentwise. For a nonconstant map the algebraic image is a curve. Away from its
finitely many singular points,
\namedref{Lemma}{lem:tubular} restricts \(u\) to the smooth image.
\namedref{Proposition}{prop:curves} supplies smooth positive approximations,
whose pullbacks prove positivity on the source outside finitely many
closed points.

Continuity removes each such point \(P\). A sufficiently small
neighbourhood is a disc over \(K(P)\), viewed over \(K\). Choose
\(0\leq w_N\leq1\), equal to one near \(P\), with supports shrinking
to \(P\), by translating one smooth profile in \(\log|t|\).
Then
\[
 \int\ddc w_N=0,\qquad \sup_N\|\ddc w_N\|_{\TV}<\infty.
\]
The bound contains only the fixed structural degree \([K(P):K]\).
A smooth scalar \(\varphi\) is constant near \(P\), since each of
its finitely many analytic unit norms is constant there. For large \(N\),
\[
 |\langle\ddc[u],\varphi w_N\rangle|
 \leq|\varphi(P)|\,
 \sup_{\supp\ddc w_N}|u-u(P)|\,\|\ddc w_N\|_{\TV}
 \longrightarrow0.
\]
Apply positivity off \(P\) to \(\varphi(1-w_N)\).
Constant curve maps pull back \(u\) to a constant.

For the extension, choose affine coordinates \((z,y)\) of an
embedding \(X\subset\Aff_K^{d+r}\), with \(z:X\to\Aff_K^d\) \'etale.
The algebraic morphism
\[
 \Theta:X\times\Aff_K^r\longrightarrow\Aff_K^{d+r},\qquad
 (a,h)\longmapsto(z(a),y(a)+h)
\]
is \'etale: after translating the second coordinate it is a base
change of \(z\). At \((x,0)\), the map induces an isomorphism of completed residue
fields. The same germ theorem \cite[Theorem 3.4.1]{Ber93} gives an
analytic inverse branch
\[
 \iota(w)=(\rho(w),y(w)-y(\rho(w)))
\]
on one ordinary ambient open. This branch agrees with the diagonal
on \(X\). Shrinking so that \(\rho\) takes values in \(O\), set
\(v=u\circ\rho\).
For any smooth algebraic curve mapping to the ambient affine space,
the fibre product with \(\Theta\) is a smooth algebraic curve.
Its map to \(X\) is covered by the first part, and the inverse
branch identifies the relevant ordinary curve open with an open
of that fibre product. Thus \(v\) is curve-positive over \(K\).
\end{proof}

\section{Curve approximation and contact rigidity}\label{sec:contact}

We first approximate the mixed currents of approachable metrics by
algebraic curves. The coefficients are chosen over the original field
\(K\), and the normalized intersection degrees remain fixed. These
degrees supply the variation bounds needed to pass the curve Green
identities to a higher-dimensional mixed current.

\subsection{Complete-intersection approximation}
\label{subsec:curve-approximation}
\begin{proposition}[Complete-intersection approximation]
\label{prop:curve-approximation}
Let \(Y/K\) be smooth projective of dimension \(n\geq2\). For
\(1\leq i<n\), let \(p_i\) be a uniform FS limit on an ample rational
line bundle \(L_i\). Then
\[
 T=\bigwedge_{i=1}^{n-1}c_1(L_i,p_i)
\]
is a weak limit, on all compact smooth tests, of individual normalized
smooth complete-intersection curve cycles \(T_a\) defined over \(K\).
For every line bundle \(H\), each \(T_a\) has normalized \(H\)-degree
\(\deg_K(HL_1\cdots L_{n-1})\).
\end{proposition}

\begin{proof}
\emph{Step 1: a smooth cut with prescribed coefficient valuations.}
Let \(Z\subset Y\) be smooth projective of pure dimension \(d\geq1\).
First consider an ample line bundle \(M\) generated by
\(s_0,\ldots,s_b\), with FS weight \(q=\max_i\log|s_i|\).
For every \(N\geq1\), choose
\begin{equation}\label{eq:deterministic-section}
 a_{i,N}\in\pi^i(1+\pi R),\qquad
 \sigma_N=\sum_{i=0}^b a_{i,N}s_i^N
\end{equation}
so that \(D_N=\divisor_Z(\sigma_N)\) is smooth.
Such a choice exists: the pure powers \(s_i^N\) have no common
zero, and characteristic-zero Bertini gives a nonempty Zariski-open
smooth locus in their coefficient space. The product
\(\prod_i\pi^i(1+\pi R)\) is Zariski dense, since a nonzero polynomial
cannot vanish on a product of infinite subsets of \(K\).
If \(Z\) has several components, impose the finitely many smoothness
conditions simultaneously. These choices are made independently
of all test forms.

\emph{Step 2: ties have zero mass for each smooth test.}
Fix a smooth \((d-1,d-1)\)-test \(\eta\) on \(Z^\an\), and let
\(\nu_\eta\) be the signed measure defined by \(\ddc\eta\).
On each component, discard the generators that vanish identically,
retaining the original indices \(i\).
By \cite[Lemma 3.2.5]{CLD12}, the compact support of \(\ddc\eta\)
avoids the zero divisors of the remaining generators.
Cover it by finitely many compact moment charts, refined by the
unit ratios \(s_i/s_j\). On each cell of the resulting finite
polyhedral complexes, the functions
\(\ell_{ij}=\log|s_i/s_j|\) are affine.

The absolute-value tropical integration formula
\cite[Section 3.7 and Proposition 3.7.2]{CLD12}
implies that \(|\nu_\eta|\) gives zero mass to inverse images of
proper affine hyperplanes in the \(d\)-dimensional cells;
lower-dimensional cells carry no top-degree mass.
Equivalently, on each compact chart one can dominate both signs
of its Lagerberg representative by a positive top form and apply
\cite[Proposition 11.1(3)]{GJRforms}. This is a statement about
total variation, not just about the signed pushforward.

Since \(|a_{i,N}|=|\pi|^i\), two summands in
\eqref{eq:deterministic-section} have the same norm exactly when
\[
 N\ell_{ij}=(i-j)\lambda.
\]
If \(\ell_{ij}\) is nonconstant on a top-dimensional cell, this
is a null hyperplane. If it is the constant \(c\), the equality
is impossible for \(c=0\) and holds for at most one positive
integer \(N\) when \(c\ne0\).
There are finitely many cells and pairs. Thus, for all sufficiently
large \(N\), the largest summand is unique at
\(|\nu_\eta|\)-almost every point.

\emph{Step 3: the one-cut estimate.}
Put \(w_N=Nq-\log|\sigma_N|\geq0\).
Unique dominance in the ultrametric inequality gives
\[
 \log|\sigma_N|
   =\max_i\{N\log|s_i|-i\lambda\},\qquad
 0\leq w_N\leq b\lambda
 \quad |\nu_\eta|\text{-almost everywhere}
\]
for all sufficiently large \(N\).
Poincar\'e--Lelong therefore gives
% Preserve the old source anchor although the argument is now deterministic.
\begin{equation}\label{eq:haar-error}
 \begin{split}
 \left|\left\langle
 \frac{[D_N]}N-c_1(M,q)\wedge[Z],\eta
 \right\rangle\right|
 &=\frac1N\left|\int_{Z^\an}w_N\,\ddc\eta\right|\\
 &\leq\frac{b\lambda}{N}\|\ddc\eta\|_{\TV}.
 \end{split}
\end{equation}
Here the threshold in \(N\) depends on \(\eta\), but the chosen
sequence \((\sigma_N)\) does not. The formula is interpreted
intrinsically on \(Z\), or after pushforward to \(Y\).
It proves one-cut convergence on every original smooth test.
The bound is a statement about integration against \(\ddc\eta\):
\(w_N\) is unbounded along \(D_N\). The argument does not give
the same estimate against an arbitrary positive background current,
whose coefficient measures may charge the tie loci.

\emph{Step 4: successive cuts and uniform limits.}
For pure-FS metrics \(p_i\), let the generating tuple lie in
\(d_iL_i\), and put \(\alpha_i=c_1(L_i,p_i)\).
At the \(i\)-th cut, restrict that tuple to the preceding smooth
intersection and apply Steps 1--3. Its divisor has class
\(N_id_iL_i\), and its current is divided by \(N_id_i\).
The Bertini conditions can also be imposed on the ambient divisors.

For clarity, let \(\mathcal C_k\) be the currents of these smooth
codimension-\(k\) complete intersections, normalized by
\(\prod_{i=1}^kN_id_i\), with \(\mathcal C_0=\{[Y]\}\).
The one-cut limit places
\(\alpha_k\wedge S\) in the weak closure of \(\mathcal C_k\)
for every \(S\in\mathcal C_{k-1}\).
Continuity in the positive background
(\namedref{Proposition}{prop:fixed-potential}) then gives inductively
\[
 \alpha_1\wedge\cdots\wedge\alpha_k
 \in\overline{\mathcal C_k}^{\,\mathrm{weak}}.
\]
Indexing a weak neighbourhood by finitely many smooth tests and
an error tolerance gives the asserted net of individual curves.
Each member has normalized \(H\)-degree
\[
 \frac{\deg_K\!\left(H\prod_{i=1}^{n-1}(N_id_iL_i)\right)}
      {\prod_{i=1}^{n-1}N_id_i}
 =\deg_K(HL_1\cdots L_{n-1}).
\]
Finally, approximate the \(p_i\) uniformly by pure-FS metrics and
take the union of the corresponding families \(\mathcal C_{n-1}\)
over all approximating tuples and tensor degrees. Every member
is still an individual normalized smooth complete intersection
with the displayed degree. Joint continuity in
\namedref{Proposition}{prop:fixed-potential} puts the limiting mixed
current in the weak closure of this union.
\end{proof}

\subsection{Concentration on a divisor}\label{subsec:divisor-limit}
The next limit uses only continuity in the positive background. In
particular, the other metrics remain fixed while the truncated
section metric tends to its divisor current.

\begin{lemma}[Concentration on a smooth divisor]\label{lem:divisor-limit}
Let \(Y/K\) be smooth projective of dimension \(n\geq1\), and let
\(s\) define a smooth effective Cartier divisor
\(i:D\hookrightarrow Y\) in a line bundle \(M\).
Let \(q,p_1,\ldots,p_{n-1}\) be globally smoothly approximable
metrics on \(M,L_1,\ldots,L_{n-1}\), respectively. Then
\begin{equation}\label{eq:divisor-limit}
 \begin{split}
 &c_1\!\left(M,\max\{q,\log|s|+N\}\right)
   \wedge\bigwedge_{j=1}^{n-1}c_1(L_j,p_j)\\
 &\hspace{2em}\longrightarrow
 i_*\!\left(\bigwedge_{j=1}^{n-1}
              c_1(L_j|_D,p_j|_D)\right)
 \end{split}
\end{equation}
as \(N\to\infty\), on every continuous scalar test on \(Y^\an\).
\end{lemma}

\begin{proof}
Put \(\rho_N=\max\{q,\log|s|+N\}\), which is smoothly approachable by
\namedref{Lemma}{lem:smooth-fs}. In a local frame,
\(\rho_N-N=\max\{q-N,\log|s|\}\searrow\log|s|\).
Local integrability and Poincar\'e--Lelong give
\(c_1(M,\rho_N)\to[D]\). Apply the fixed-potential continuity of
\namedref{Proposition}{prop:fixed-potential} in the remaining slots.
The projection formula for approachable products
\cite[Proposition 8.7.5]{CLD25} identifies the limit with the
right-hand side of \eqref{eq:divisor-limit}.
All measures are positive and have the common mass
\(\deg_K(ML_1\cdots L_{n-1})\), by \namedref{Lemma}{lem:degree}.
Uniform smooth density extends convergence to continuous tests.
\end{proof}

\subsection{Contact rigidity}\label{subsec:contact-proof}
The following criterion compares two metrics on the same ample line.
Only the lower metric is assumed to be a uniform FS limit; the upper
metric is tested on algebraic curves. The proof uses the preceding
curve approximation and divisor limit, independently of envelope continuity.

\begin{theorem}[Contact rigidity]\label{thm:contact}
Let \(Y/K\) be smooth, projective, and integral of dimension
\(n\geq1\), and let \(L\) be ample. Let \(p\) be a uniform FS
limit on \(L\), and let \(r\) be a finite continuous metric on
\(L\) whose pullback by every morphism from a smooth algebraic
\(K\)-curve has positive Chern current. If
\[
 g:=r-p\geq0,
 \qquad
 \int_{Y^\an}g\,c_1(L,p)^n=0,
\]
then \(r=p\).
\end{theorem}

\begin{proof}
\emph{Step 1: pass from curves to mixed measures.}
Put
\[
 \alpha=c_1(L,p),\qquad T=\alpha^{n-1},\qquad
 \mu=\alpha^n,\qquad D=\deg_K(L^n)>0.
\]
For \(n>1\), approximate \(T\) by the normalized smooth curve
cycles \(T_a\) of \namedref{Proposition}{prop:curve-approximation}; for
\(n=1\), take \(T_a=[Y]\). The pullbacks of \(r\) and \(p\)
to each curve have positive Chern measures. Their normalized
pushforwards \(\sigma_a,\mu_a\) have mass \(D\), and
\[
 \nu_a=\ddc(gT_a)=\sigma_a-\mu_a,
 \qquad\|\nu_a\|_{\TV}\leq2D.
\]
By \namedref{Proposition}{prop:fixed-potential}, \(\mu_a\to\mu\).
\namedref{Lemma}{lem:continuous-coefficients} gives
\(\nu_a\to\nu:=\ddc(gT)\) on smooth tests. The variation bounds
extend these limits to continuous scalar tests. Thus
\(\sigma_a\to\sigma:=\nu+\mu\), where \(\sigma\) is a positive
measure of mass \(D\).

\emph{Step 2: use zero energy to annihilate the mixed Laplacian.}
We pass the curve Green identities of \namedref{Lemma}{lem:curve-green}
to the limit. For a fixed smooth function \(h\),
\namedref{Lemma}{lem:continuous-coefficients} bounds the total variations of
\(\ddc h\wedge T_a\): a finite collection of local positive test
evaluations bounds the fixed coefficient form \(\ddc h\).
Consequently,
\begin{equation}\label{eq:contact-energy}
 \int g\,\ddc(hT)=\int h\,d\nu,
 \qquad
 -\int(g+th)\ddc((g+th)T)\geq0\quad(t\in\R).
\end{equation}
All integrals are over \(Y^\an\). The contact assumption implies
\[
 0\leq-\int g\,d\nu
   =\int g\,d\mu-\int g\,d\sigma\leq0.
\]
The quadratic polynomial in \(t\) from
\eqref{eq:contact-energy} therefore has zero constant term.
Green symmetry writes its linear term as \(-2t\int h\,d\nu\).
Nonnegativity for every \(t\in\R\) forces that coefficient to
vanish. Since this holds for every smooth \(h\), we have
\(\nu=0\).

If \(\rho\) is any finite continuous curve-positive metric on
\(mL\), set \(b=\rho-mp\). Its normalized curve Laplacians have
variations at most \(2mD\). The same passage of Green symmetry,
now for the two continuous gaps \(g\) and \(b\), gives
\begin{equation}\label{eq:contact-green}
 \int g\,\ddc(bT)=\int b\,d\nu=0.
\end{equation}
Only products with the approachable lower metric \(p\) have been
used; the upper metric \(r\) enters through its curve restrictions.

\emph{Step 3: restrict the contact identity and argue by dimension.}
When \(n=1\), \namedref{Lemma}{lem:curve-green} and the zero energy
already proved make \(g\) constant. Its integral against the
positive measure \(\mu\) of mass \(D\) makes that constant zero.
Suppose \(n>1\) and proceed by induction. Fix a closed point
\(x\in Y\). For sufficiently large \(m\), choose a smooth
\(H=\divisor(s)\) in \(mL\) through \(x\). To impose this
condition, use \(H^0(Y,mL\otimes\mathcal I_x)\): for \(m\gg0\)
it is generated away from \(x\) and realizes arbitrary first jets
at \(x\). A nonzero first jet ensures smoothness there, and
characteristic-zero Bertini ensures smoothness elsewhere. The
nonempty open conditions admit a \(K\)-rational choice since
\(K\) is infinite.

Set \(\rho_N=\max\{mp,\log|s|+N\}\). These metrics are smoothly
approximable by \namedref{Lemma}{lem:smooth-fs}, hence curve-positive.
Equation~\eqref{eq:contact-green}, together with
\(\int g\,d\mu=0\), yields
\[
 \int_{Y^\an}g\,c_1(mL,\rho_N)\wedge\alpha^{n-1}=0.
\]
\namedref{Lemma}{lem:divisor-limit} now gives
\[
 \int_{H^\an}g|_H\,c_1(L|_H,p|_H)^{n-1}=0.
\]
The integrand is nonnegative. It therefore has zero integral on
each smooth integral component of \(H\). The curve condition
passes to that component by composition, \(L|_H\) is ample, and
\(p|_H\) is a uniform FS limit by restriction of the generating
tuples. Induction gives \(g(x)=0\). Closed points give a dense
set of rigid points of \(Y^\an\); continuity implies \(g=0\)
everywhere.
\end{proof}

\section{Envelopes on a fixed regular model}\label{sec:envelopes}

For an ample line bundle \(L\) on a projective \(K\)-variety and a
continuous metric \(r\), set
\begin{equation}\label{eq:envelope-definition}
 P_L(r)=\sup\{q\le r:q\text{ is a continuous semipositive metric on }L\}.
\end{equation}
We use the standard finiteness, translation, scaling, and monotonicity
properties of this envelope \cite[Remark~6.2 and Proposition~6.3]{BGM}.
In particular,
\begin{equation}\label{eq:envelope-lipschitz}
 \norm{P_L(r)-P_L(r')}_\infty\le \norm{r-r'}_\infty.
\end{equation}
The envelope is not assumed continuous at this stage.

\subsection{Test ideals of admissible section divisors}\label{subsec:test-ideals}

Let \(\cX/R\) be regular, integral, projective and flat, with smooth
generic fibre \(Y\) of dimension \(n\ge1\). Extend \(L\) to a line bundle
\(\cL\) on \(\cX\), and write \(\theta\) for its model metric.
Such an extension exists because the closure of a generic Cartier divisor
is Cartier on the regular scheme \(\cX\).
After clearing a fixed tensor denominator and adjusting the reference
identification by a power of \(\pi\), a model metric \(r\) can be written
\begin{equation}\label{eq:model-obstacle}
 r=\theta+\log|\cE|,
 \qquad \cE=\OO_{\cW}(-F)\subseteq\OO_{\cW},
\end{equation}
where \(f:\cW\to\cX\) is a normal projective model inducing the identity
on \(Y\), and \(F\) is an effective vertical Cartier divisor.
Put \(p=P_L(r)\).

\begin{proposition}[The test-ideal estimate]\label{prop:test-ideal-estimate}
There are coherent vertical ideals \(\mathfrak j_m\subseteq\OO_{\cX}\),
for \(m\ge1\), and an integer \(B\ge0\), independent of \(m\), such that
\begin{equation}\label{eq:test-ideal-bounds}
 \pi^B\mathfrak j_m\OO_{\cW}\subseteq\cE^m,\qquad
 p\le\theta+\psi_m\le r+\frac{B\lambda}{m},
 \qquad \psi_m=\frac1m\log|\mathfrak j_m|.
\end{equation}
For every globally generated ample Cartier divisor \(H\) on \(\cX\),
\begin{equation}\label{eq:test-ideal-generation}
 \mathfrak j_m\cL^{\otimes m}\otimes
 \OO_{\cX}\bigl(K_{\cX}+(n+1)H\bigr)
 \quad\text{is globally generated.}
\end{equation}
\end{proposition}

\begin{proof}
For \(l\ge1\), let
\[
 V_l=H^0\bigl(\cW,f^*\cL^{\otimes l}\otimes\cE^l\bigr).
\]
Since \(f_*\OO_{\cW}=\OO_{\cX}\), every member of \(V_l\) is a section
of \(\cL^{\otimes l}\) on \(\cX\). Normality also gives
\begin{equation}\label{eq:weighted-unit-ball}
 V_l=\{s\in H^0(Y,lL):\norm{s}_{lr,\sup}\le1\}.
\end{equation}
Indeed, a generic section has no horizontal poles, while the norm
inequalities at vertical divisorial points are exactly the remaining
codimension-one extension conditions on \(\cW\).
For \(0\ne s\in V_l\), denote its effective divisor on \(\cX\) by \(D_s\).
Then
\[
 D_s\sim l\cL,\qquad f^*D_s\ge lF.
\]
Define
\begin{equation}\label{eq:divisor-sum}
 \mathfrak j_m
 =\sum_{l\ge1}\ \sum_{0\ne s\in V_l}
       \tau_{\cX}\!\left(\frac mlD_s\right).
\end{equation}
This sum equals a finite subsum by Noetherianity.
For a sufficiently large \(l\), a generating tuple of \(H^0(Y,lL)\)
can be scaled into \(V_l\). At every point of \(Y\) one of these
sections is nonvanishing. Generic-fibre comparison with multiplier
ideals \cite[Corollary~7.21(e)]{BMP} therefore shows that
\(\mathfrak j_m|_Y=\OO_Y\). Thus \(\mathfrak j_m\) is vertical.

\smallskip\noindent
\emph{Lower bound.}
Write \(\tau_0=\tau_{\cX}(0)\). This fixed ideal contains a power of \(\pi\),
since its generic restriction is \(\OO_Y\).
Cartier translation and subadditivity on the regular scheme \(\cX\)
give, for every \(k\ge1\),
\begin{equation}\label{eq:lower-ideal}
 \OO_{\cX}(-kmD_s)\tau_0
 =\tau_{\cX}(kmD_s)
 \subseteq
 \tau_{\cX}\!\left(\frac mlD_s\right)^{kl}
 \subseteq\mathfrak j_m^{kl}.
\end{equation}
Here we use \cite[Proposition~7.16(b) and Theorem~7.22(e)]{BMP}.
Taking logarithmic norms, dividing by \(kml\), and letting \(k\to\infty\)
yields
\[
 \frac1l\log\norm{s}_{l\theta}\le\psi_m.
\]
If a pure-FS metric \(q\) lies below \(r\), its defining sections belong
to the corresponding \(V_l\) by \eqref{eq:weighted-unit-ball}; taking their
maximum gives \(q\le\theta+\psi_m\).
Every continuous semipositive competitor below \(r\) is a uniform
limit of pure-FS metrics shifted down by constants in
\(\lambda\mathbf Q\) tending to zero
\cite[Corollary~7.6 and Theorem~7.8]{BE}.
The lower bound in \eqref{eq:test-ideal-bounds} follows.

\smallskip\noindent
\emph{Upper bound.}
Choose compatible canonical divisors and set
\(\Delta_{\cW}=f^*K_{\cX}-K_{\cW}\).
The fractional test ideal \(\tau_{\cW}(\Delta_{\cW})\) is defined because
\(K_{\cW}+\Delta_{\cW}\) is Cartier
\cite[Definition~7.20]{BMP}.
Birational inclusion, monotonicity, and Cartier translation give
\begin{align}
 \tau_{\cX}\!\left(\frac mlD_s\right)\OO_{\cW}
 &\subseteq
 \tau_{\cW}\!\left(\Delta_{\cW}+\frac ml f^*D_s\right) \notag\\
 &\subseteq
 \tau_{\cW}(\Delta_{\cW}+mF)
 =\cE^m\tau_{\cW}(\Delta_{\cW}).
 \label{eq:birational-bound}
\end{align}
The cited properties are
\cite[Proposition~7.16(a),(b) and Theorem~7.22(d)]{BMP}.
The fixed fractional ideal on the right is coherent and restricts to
\(\OO_Y\). A single power \(\pi^B\) clears its denominators on a finite
affine cover of \(\cW\). Summing \eqref{eq:birational-bound} gives
\(\pi^B\mathfrak j_m\OO_{\cW}\subseteq\cE^m\), and hence the upper bound.

\smallskip\noindent
\emph{Global generation.}
Each effective boundary \(\Delta=(m/l)D_s\) is
\(\mathbf Q\)-linearly equivalent to \(m\cL\).
For \(M=m\cL+K_{\cX}+H\), one has
\(M-K_{\cX}-\Delta\sim_{\mathbf Q}H\).
Since the special fibre has dimension \(n\),
\cite[Corollary~7.24]{BMP} makes
\(\tau_{\cX}(\Delta)\cL^{\otimes m}(K_{\cX}+(n+1)H)\)
globally generated. A finite subsum in \eqref{eq:divisor-sum}
proves \eqref{eq:test-ideal-generation}.
\end{proof}

\subsection{Uniform approximation and orthogonality}\label{subsec:envelope-proof}

\begin{theorem}[Envelopes on a regular model]\label{thm:envelope}
Let \(\cX/R\) be regular, integral, projective and flat, with smooth generic
fibre \(Y\) of dimension \(n\), and let \(L\) be an ample rational line
bundle on \(Y\).
For every finite continuous metric \(r\) on \(L\), the envelope
\(p=P_L(r)\) is continuous, is a uniform limit of pure-FS metrics, and
satisfies
\begin{equation}\label{eq:regular-model-orthogonality}
 \int_{Y^{\an}}(r-p)c_1(L,p)^n=0.
\end{equation}
Here the measure is the CLD product.
For a model obstacle, after clearing fixed tensor denominators,
there are pure-FS metrics \(q_\ell\), for all sufficiently large
integers \(\ell\), such that
\begin{equation}\label{eq:regular-model-approximation}
 0\le q_\ell-p\le \frac C\ell,
\end{equation}
where \(C\) depends on the obstacle and the fixed auxiliary data.
\end{theorem}

\begin{proof}
If \(n=0\), then \(Y=\operatorname{Spec}F\) for a finite extension
\(F/K\), and \(Y^\an\) is a point. In a fixed frame, weights of pure-FS
metrics include \(\lambda\mathbf Q\), so every continuous metric is
semipositive and \(p=r\). A model metric is pure-FS after clearing its
fixed tensor denominator, which also gives the stated approximation.
Assume henceforth that \(n\ge1\).

First treat a model obstacle in the form \eqref{eq:model-obstacle}.
Choose a globally generated ample divisor \(H\) sufficiently positive that
\[
 \cA=\OO_{\cX}\bigl(K_{\cX}+(n+1)H\bigr)
\]
is ample and globally generated. Put \(A=\cA|_Y\), and let \(q_A\)
be its model FS metric.
Choose \(k\ge1\) such that \(kL-A\) is very ample and a pure-FS metric
\(\nu\) defined by a generating tuple on that line bundle.
Multiplying this tuple by a sufficiently large negative power of \(\pi\),
arrange that
\[
 h:=\frac{q_A+\nu}{k}\ge r.
\]
By \namedref{Proposition}{prop:test-ideal-estimate}, the generators of
\(\mathfrak j_m\cL^{\otimes m}\cA\) define a pure-FS metric
\[
 t_m=m\theta+q_A+\log|\mathfrak j_m|
 \quad\text{on }mL+A.
\]
Tensoring with the tuple defining \(\nu\) gives a pure-FS metric on \(L\),
\[
 h_m=\frac{t_m+\nu}{m+k}
     =\frac m{m+k}(\theta+\psi_m)+\frac k{m+k}h.
\]
Consequently,
\[
 p\le h_m\le r+\frac C{m+k},
 \qquad C=B\lambda+k\norm{h-r}_\infty.
\]
The semipositive metric \(h_m-C/(m+k)\) lies below \(r\), hence below
\(p\). Thus
\[
 0\le h_m-p\le\frac C{m+k}.
\]
These metrics are represented by generating sections in degree \(m+k\).
Relabeling \(\ell=m+k\) proves \eqref{eq:regular-model-approximation}.
A fixed tensor denominator is divided out using
\(P_{dL}(dr)=dP_L(r)\).

For a continuous obstacle, use uniform density of model metrics
\cite[Proposition~5.16(vi) and Theorem~5.18]{BE}.
Equation~\eqref{eq:envelope-lipschitz} shows that their envelopes converge
uniformly to \(P_L(r)\), and a diagonal choice gives pure-FS approximation.
Finally, \(Y\) is geometrically reduced, so
\cite[Theorem~6.7]{BGM} gives \eqref{eq:regular-model-orthogonality}.
\namedref{Lemma}{lem:degree} identifies its measure with the CLD product.
\end{proof}

\begin{corollary}[Projective space]\label{cor:projective-envelope}
Let \(a,b\ge1\) and \(L=\OO_{\PP^a_K}(b)\).
The conclusions of \namedref{Theorem}{thm:envelope} hold for every continuous metric
on \(L\). For a model obstacle represented as in
\eqref{eq:model-obstacle} on a model above \(\PP^a_R\), the generators
of \(\mathfrak j_m\OO(bm)\) give pure-FS metrics \(q_m\) satisfying
\[
 0\le q_m-P_L(r)\le B\lambda/m.
\]
\end{corollary}

\begin{proof}
Take \(\cX=\PP^a_R\).
For the more precise assertion, choose the hyperplane divisor \(H\);
then \(K_{\cX}+(a+1)H=0\).
Equation~\eqref{eq:test-ideal-generation} makes
\(q_m=\theta+m^{-1}\log|\mathfrak j_m|\) pure-FS.
By \eqref{eq:test-ideal-bounds}, \(p\le q_m\le r+B\lambda/m\);
the envelope property applied to \(q_m-B\lambda/m\) finishes the proof.
\end{proof}

\section{Proof of local regularization}\label{sec:local-proof}
The next lemma places a local curve-positive function inside one
global metric on projective space. It uses only algebraic sections
and the maximum property on curves.

\subsection{A projective extension}\label{subsec:projective-extension}

\begin{lemma}[Projective extension]\label{lem:projective-extension}
Let \(W\subset\Aff_K^{a,\an}\) be an ordinary open, \(a\ge1\),
and let \(v:W\to\mathbf R\) be finite, continuous and curve-positive.
For every \(x\in W\), there are \(d\ge1\), a finite continuous
curve-positive metric \(r\) on \(\OO_{\PP^a_K}(d)\), a section \(s_0\)
nonvanishing near \(x\), and a constant \(c\in\mathbf R\), such that
\[
 -\log\norm{s_0}_r=v+c
 \quad\text{on an ordinary neighbourhood of }x.
\]
\end{lemma}

\begin{proof}
Set \(c=\lambda/4-v(x)\).
After shrinking \(W\), continuity gives
\begin{equation}\label{eq:normalization}
 0<v+c<\lambda/2\quad\text{on }W.
\end{equation}
Choose a compact polynomial norm box
\[
 C=\{|F_i|\le|c_i|,\ |G_j|\ge|d_j|\}\subset W
\]
with all defining inequalities strict at \(x\), including upper bounds
on the affine coordinates.
Here \(F_i,G_j\in K[z_1,\ldots,z_a]\) and \(c_i,d_j\in K^\times\).
Such boxes form a neighbourhood basis: choose thresholds in
\(|K^\times|^{\mathbf Q}\) inside an ordinary basic neighbourhood and
encode them by powers of the polynomials.

For \(N\ge1\) sufficiently large, every logarithmic term in the following
maximum is negative on a neighbourhood of \(x\):
\begin{equation}\label{eq:collar-weight}
 b_0=\max\left\{
 0,\ \log\left|\frac{F_i^N}{\pi c_i^N}\right|,
       \log\left|\frac{d_j^N}{\pi G_j^N}\right|
 \right\}
\end{equation}
On the open where every \(G_j\) is nonzero, this function
satisfies \(b_0=0\) near \(x\) and \(\{b_0\le\lambda\}=C\).

Put \(t_0=\prod_jG_j^N\), with the empty product equal to \(1\).
Clear the denominators in \eqref{eq:collar-weight} and homogenize
\[
 t_0,\qquad
 \frac{t_0F_i^N}{\pi c_i^N},\qquad
 \frac{t_0d_j^N}{\pi G_j^N}
\]
to a common degree \(d\), chosen strictly larger than \(\deg t_0\).
Write \(s_0\) for the homogenization of \(t_0\).
Add sufficiently small \(K\)-multiples of
\(Z_0^d,\ldots,Z_a^d\), so that their ratios to \(s_0\) have norm
less than \(1\) on \(C\).
The resulting tuple generates \(\OO(d)\); let \(q\) be its FS metric.
Then
\[
 b:=-\log\norm{s_0}_q\in[0,+\infty]
\]
equals \(b_0\) on \(C\), vanishes near \(x\), and satisfies
\begin{equation}\label{eq:collar-sublevel}
 \{b\le\lambda\}=C.
\end{equation}
Indeed, \(s_0\) vanishes at infinity because \(d>\deg t_0\), and the
remaining original ratios impose precisely the inequalities defining \(C\).
The added sections preserve those inequalities on \(C\).
We interpret \(b=+\infty\) on the zero divisor of \(s_0\).

The ordinary opens
\[
 U=\{b<\lambda\},\qquad V=\{b>\lambda/2\}
\]
cover \(\PP_K^{a,\an}\), and \(U\subset W\).
On \(U\), define a metric by its \(s_0\)-weight
\(\max\{v+c,b\}\); on \(V\), use \(q\).
The definitions agree on \(U\cap V\) by \eqref{eq:normalization}.
They therefore give a finite continuous metric \(r\).
Near \(x\), one has \(b=0<v+c\), so its weight is \(v+c\).

On every smooth algebraic \(K\)-curve, the weights just used are
continuous subharmonic functions or maxima of two such functions.
Their maxima are subharmonic by the curve regularization of
\namedref{Proposition}{prop:curves} and the regularized maximum construction.
Thus \(r\) is curve-positive.
\end{proof}

\subsection{Completion of the proof}\label{subsec:regularization-proof}

\begin{proof}[Proof of \namedref{Theorem}{thm:main}]
Local uniform convergence of smooth positive potentials implies current
positivity by testing against \(\ddc\eta\).
Conversely, fix \(x\in O\) on a positive-dimensional component of \(X\).
\namedref{Proposition}{prop:curve-restriction} supplies an ordinary ambient open
\(W\subset\Aff_K^{a,\an}\) and a finite continuous curve-positive
function \(v\) agreeing with \(u\) along \(X^{\an}\cap W\).
Apply \namedref{Lemma}{lem:projective-extension}, and retain its \(r,s_0,c\).

By \namedref{Corollary}{cor:projective-envelope}, the envelope \(p=P_{\OO(d)}(r)\) is a uniform
FS limit and
\[
 p\le r,\qquad
 \int_{\PP_K^{a,\an}}(r-p)c_1(\OO(d),p)^a=0.
\]
\namedref{Theorem}{thm:contact} gives \(p=r\).
By \namedref{Lemma}{lem:smooth-fs}, choose global smooth positive metrics \(r_j\) with
\(\norm{r_j-r}_\infty\le4^{-j}\).
On one fixed ordinary neighbourhood \(V_x\) of \(x\) in \(X^{\an}\)
where the weight identity holds, put
\[
 u_j=-\log\norm{s_0}_{r_j}\big|_{V_x}-c.
\]
These functions are smooth, have positive curvature, and satisfy
\(\norm{u_j-u}_{V_x}\le4^{-j}\).
For the decreasing assertion, define
\[
 w_j=u_j+2\cdot4^{-j}.
\]
Then
\[
 u+4^{-j}\le w_j\le u+3\cdot4^{-j},
 \qquad w_{j+1}\le u+\tfrac34\,4^{-j}<w_j.
\]
Thus \(w_j\downarrow u\) uniformly on the same \(V_x\).
Constant translations preserve curvature.
On a zero-dimensional component, every finite continuous function is
locally constant and already smooth.
\end{proof}

\section{Consequences for continuous potentials}\label{sec:consequences}

We first apply local regularization to the operation constructed in
\namedref{Proposition}{prop:fixed-potential}. The background current remains
arbitrary; in particular, it need not itself be a product of curvature
currents.

\subsection{Proof of the mixed-current theorem}
\label{subsec:mixed-current-proof}
\begin{proof}[Proof of \namedref{Theorem}{thm:background-products}]
\namedref{Theorem}{thm:main} places every \(u_i\) in \(\Ploc(O)\).
Iterating \namedref{Proposition}{prop:fixed-potential} in
\eqref{eq:recursive-product} proves positivity, closedness, symmetry,
and joint continuity, including nets. Locality, additivity, and
positive homogeneity follow from the same recursion. A positive
current of top degree is a Radon measure.

It remains to prove independence of the potentials. Let \(h\) be
finite continuous with \(\ddc[h]=0\), and let \(T\) be any positive
closed current for which the bidegrees permit the operation. Both
\(h\) and \(-h\) have positive curvature currents, so
\namedref{Theorem}{thm:main} and \namedref{Proposition}{prop:fixed-potential} give
\[
 \ddc(hT)\geq0,\qquad
 -\ddc(hT)=\ddc((-h)T)\geq0.
\]
The positive cone is pointed by \namedref{Lemma}{lem:continuous-coefficients}; hence
\begin{equation}\label{eq:continuous-kernel}
 \ddc[h]=0\quad\Longrightarrow\quad\ddc(hT)=0.
\end{equation}
If \(\ddc[w]=\ddc[u_i]\), apply this implication to \(h=w-u_i\)
and to the positive closed product of the other slots. Symmetry permits
us to place the changed slot first, proving the assertion.
\end{proof}

\subsection{Chern products and closure operations}
\begin{corollary}[Continuous Chern products]
\label{cor:continuous-chern-products}
Let \(X/K\) be smooth quasi-projective of pure dimension \(n\), let
\(O\subset X^\an\) be an ordinary open, and let \(S\) be a positive
closed CLD current of bidegree \((p,p)\) on \(O\). Let
\(\overline L_1,\ldots,\overline L_m\) be continuously metrized
analytic line bundles on \(O\), with finite local weights and positive
CLD curvature currents, where \(m\leq n-p\). Their local products
glue to a positive closed current
\[
 c_1(\overline L_1)\wedge\cdots\wedge c_1(\overline L_m)\wedge S.
\]
It depends only on the curvature currents and \(S\), agrees with the
smooth wedge product when all weights are smooth, and is a positive
Radon measure if \(m+p=n\). The same assertions hold for rational line
bundles after clearing and dividing their tensor denominators.
\end{corollary}

\begin{proof}
Apply \namedref{Theorem}{thm:background-products} in local frames. On overlaps,
weights differ by logarithms of analytic units, whose curvature is zero
by Poincar\'e--Lelong. Equation~\eqref{eq:continuous-kernel} identifies
the products. The remaining assertions, including rational
normalization, follow from the same theorem.
\end{proof}

\begin{corollary}[Products with a smooth subvariety]
\label{cor:smooth-subvariety-products}
Let \(i:Z\hookrightarrow X\) be a closed immersion of smooth
quasi-projective \(K\)-varieties, with \(X\) of pure dimension \(n\)
and \(Z\) of pure dimension \(d\).
Let \(O\subset X^\an\) be ordinary open, put
\(O_Z=(i^\an)^{-1}(O)\), and let \(u_1,\ldots,u_m\) be finite
continuous functions on \(O\) with \(\ddc[u_j]\geq0\), where
\(m\leq d\). Write \([Z]_O=i_*[O_Z]\) for the integration current
on \(O\), and write \(\widetilde u_j=u_j|_{O_Z}\). Then
\begin{equation}\label{eq:subvariety-projection}
 \mathcal B_{[Z]_O}(u_1,\ldots,u_m)
 =i_*\mathcal B_{[O_Z]}(\widetilde u_1,\ldots,\widetilde u_m).
\end{equation}
\end{corollary}

\begin{proof}
Positive smooth approximants from \namedref{Theorem}{thm:main} restrict to
positive smooth approximants on \(O_Z\). The assertion is then the
projection formula for approachable mixed products
\cite[Proposition 8.7.5]{CLD25}, applied to the closed immersion \(i\).
\end{proof}

Formula~\eqref{eq:subvariety-projection} identifies a product with the
integration current of \(Z\) with the intrinsic product on \(Z\).
It does not require a pullback operation for arbitrary currents.

\begin{corollary}[Closure operations]\label{cor:closure-operations}
Let \(u_1,\ldots,u_r\) be finite continuous current-positive functions
on an ordinary open \(O\subset X^\an\), where \(X/K\) is smooth
quasi-projective. If \(F:\R^r\to\R\) is convex and nondecreasing in
each coordinate, then \(F(u_1,\ldots,u_r)\) is current-positive.
In particular, finite maxima and sums with nonnegative coefficients
preserve this class.

If \(V\subset Y^\an\) is an ordinary open in a smooth
quasi-projective \(K\)-variety and \(f:V\to O\) is an analytic
morphism, then \(u_i\circ f\) is current-positive on \(V\).
\end{corollary}

\begin{proof}
By \namedref{Theorem}{thm:main}, the \(u_i\) are locally uniformly
approachable by smooth positive potentials. The convex composition
assertion is therefore \cite[Corollary 8.2.4]{CLD25}.
For the pullback assertion, pull back those smooth approximants;
smooth pullback preserves positivity, and locally uniform convergence
is preserved because continuous maps send compact sets to compact sets.
\end{proof}

\subsection{The projective curve criterion}
\begin{corollary}[A projective curve criterion]
\label{cor:projective-criterion}
Let \(\mathscr Y/R\) be a regular integral projective flat model with
smooth generic fibre \(Y\) of positive dimension. Let \(L\) be an
ample line bundle on \(Y\), and let \(r\) be a finite continuous
metric on \(L\). The following conditions are equivalent:
\begin{enumerate}
\item For every algebraic morphism \(f:C\to Y\) from a smooth
\(K\)-curve, the pulled-back metric \(f^*r\) has positive Chern current.
\item The original CLD current \(c_1(L,r)\) is positive.
\item The metric \(r\) is semipositive, equivalently a uniform limit
of pure Fubini--Study metrics on \(L\).
\item The metric \(r\) is a uniform limit of global smooth metrics on
\(L\) with positive curvature forms.
\end{enumerate}
The criterion also holds for an ample rational line bundle by clearing
its tensor denominator.
\end{corollary}

\begin{proof}
Suppose the first condition holds. \namedref{Theorem}{thm:envelope} gives a
finite continuous semipositive envelope \(p=P_L(r)\), which is a
uniform Fubini--Study limit and satisfies
\[
 p\leq r,\qquad
 \int_{Y^\an}(r-p)c_1(L,p)^{\dim Y}=0.
\]
\namedref{Theorem}{thm:contact} applies to the curve-positive upper metric
\(r\) and gives \(r=p\). This proves the third condition.
\namedref{Lemma}{lem:smooth-fs} gives the fourth condition from the third,
and positivity passes to the current limit, giving the second.
Finally, \namedref{Proposition}{prop:curve-restriction}, applied in local
frames, gives the first condition from the second. Positive tensor
powers scale each curvature by the same positive integer, which proves
the rational-line assertion.
\end{proof}

\section{Comparison with adelic mixed measures}\label{sec:adelic}

\begin{definition}[Boundary completion]\label{def:boundary}
Let \(U/k\) be quasi-projective. We use the rational boundary
completion of \cite[Sections 2.4--2.5]{YZ}, retaining its specified
line-bundle identifications. A boundary datum consists of a
projective compactification, an effective Cartier boundary with
its multiplicities, and strictly positive Green data.
On common dominating models a Cauchy presentation satisfies, for
every \(\varepsilon>0\) and all sufficiently large \(j,l\),
\begin{equation}\label{eq:boundary-order}
 -\varepsilon\overline\Delta_\partial
 \leq\overline D_j-\overline D_l
 \leq\varepsilon\overline\Delta_\partial .
\end{equation}
In a fixed local identification at \(v\), this implies
\[
 |w_j-w_l|\leq\varepsilon b_v\quad\text{on }U_K^\an,
\]
where the boundary Green function \(b_v\) is finite continuous
on that open. The completion uses the full order
\eqref{eq:boundary-order}, not merely this scalar inequality.
A strongly nef adelic line bundle has a Cauchy presentation by
nef model objects. An integrable one is a difference of two
strongly nef objects, allowing rational tensor powers and the
specified isometries \cite[Definition 2.5.2]{YZ}.
\end{definition}

\subsection{Identification and positivity of the prescribed measure}
\label{subsec:adelic-proof}
\begin{proof}[Proof of \namedref{Theorem}{thm:adelic}]
Choose integrable presentations with their prescribed isometries,
\[
 \overline L_i^{\otimes r_i}
 \simeq\overline P_i\otimes\overline Q_i^{-1},
 \qquad r_i\in\Z_{>0},
\]
where \(\overline P_i,\overline Q_i\) are strongly nef. Their metrics
at \(v\) are locally psh-approachable, and their mixed measures are
the stable CLD products
\cite[Lemma 3.6.7 and Sections 3.6.7--3.6.8]{YZ}.
Let \(u_i,p_i,q_i\) be weights in compatible local frames. The
specified isometries give
\[
 r_i u_i=p_i-q_i+h_i,\qquad
 h_i=-\log|a_i|,\qquad \ddc h_i=0,
\]
for analytic units \(a_i\). By the definition in
\cite[Section 3.6.8]{YZ}, the prescribed signed measure is
\begin{equation}\label{eq:adelic-expansion}
 \mu_v^{\mathrm{YZ}}
 =\frac1{\prod_i r_i}
   \sum_{\epsilon\in\{0,1\}^n}(-1)^{|\epsilon|}
       \prod_{i=1}^n c_1(\overline R_{i,\epsilon_i})_v,
 \qquad
 \overline R_{i,0}=\overline P_i,\quad
 \overline R_{i,1}=\overline Q_i.
\end{equation}
Its independence of presentation is part of that construction.

\namedref{Theorem}{thm:main} makes each target weight \(u_i\) approachable.
Apply additivity of approachable products to
\(r_i u_i+q_i=p_i+h_i\) in every slot; the zero-curvature terms
\(h_i\) disappear by \eqref{eq:continuous-kernel}.
The alternating expansion is exactly
\begin{equation}\label{eq:adelic-local}
 \mu_v^{\mathrm{YZ}}=\bigwedge_{i=1}^n\ddc u_i.
\end{equation}
The right-hand side is positive by
\namedref{Theorem}{thm:background-products}. Positivity of Radon measures
is local, so this proves the assertion on all of \(U_K^\an\).
\end{proof}

\appendix
\section{A quantitative comparison of mixed measures}
\label{app:comparison}

For completeness, we record a comparison estimate in which the constant
is independent of the metrics being compared. It is separate from the
proof of contact rigidity: the divisor limit there uses only the
continuity of a fixed potential acting on positive closed currents.

\begin{lemma}[Global domination of finitely many forms]
\label{lem:global-calibration}
Let \(Y/K\) be smooth projective, and let
\(\eta_1,\ldots,\eta_s\) be smooth symmetric real \((1,1)\)-forms
on \(Y^\an\). There are an ample line bundle \(H\), a global smooth
metric \(\phi\) on \(H\), and its positive curvature
\(\omega=c_1(H,\phi)\), such that
\[
 -\omega\leq\eta_i\leq\omega\qquad(1\leq i\leq s).
\]
The inequalities are in the original CLD form cone, and all the
constructions take place over \(K\).
\end{lemma}

\begin{proof}
Choose a very ample line bundle \(A\) on \(Y\). We first construct
one global positive curvature that controls the square of a local
moment coordinate. On a compact polynomial norm box \(C\) contained
in an affine chart of the embedding defined by \(A\), let \(f\) be
an analytic unit on a neighbourhood of \(C\). The box is defined by
finitely many upper bounds and positive lower bounds on polynomial
functions, with bounds in \(|K^\times|^\Q\). After clearing powers,
its affinoid algebra is the completion of the affine coordinate
algebra with the lower-bounded polynomials inverted.
Thus \(f\) can be uniformly approximated on \(C\) by a rational
function \(h\) from this algebra. Choose the error smaller than
\(\min_C|f|\). The ultrametric inequality gives
\[
 |h|=|f|\quad\text{throughout }C.
\]
In particular \(h\) is a rational unit there. This equality of
moment maps gives the same differential forms on \(\Int C\),
by \cite[Lemma 3.1.10]{CLD12}.

Homogenizing a numerator and denominator for \(h\) expresses it
as \(h=s_1/s_0\), where \(s_0,s_1\) are sections of some
positive tensor power \(A^d\), both nonzero on \(C\).
Set \(M=\|\log|h|\|_C\), and choose
\(\varepsilon>2M\). In the construction of
\namedref{Proposition}{prop:regularized-max}, choose \(\rho\) strictly
positive on \((-1,1)\), and write
\[
 H_\varepsilon(t)=R_{2,\varepsilon}(t,0).
\]
Then
\[
 H_\varepsilon''(t)
 =\frac1\varepsilon
   (\rho*\check\rho)(t/\varepsilon)>0
       \qquad(|t|<2\varepsilon),
 \qquad\check\rho(t)=\rho(-t).
\]
Hence \(H_\varepsilon''\geq c>0\) on \([-M,M]\).

Take a global generating tuple of \(A^d\), multiply its members
by sufficiently small elements of \(K^\times\), and append them
to \(s_0,s_1\). If \(a\) is any appended section, choose the
coefficient so that
\[
 \log|a/s_0|<-2\varepsilon-1
 \quad\text{on }C
\]
where the value is \(-\infty\) at a zero. Thus its logarithmic
entry lies more than \(2\varepsilon\) below the maximum of
\(0\) and \(\log|h|\) throughout \(C\).
Every appended entry is exactly inactive in the regularized
maximum on a neighbourhood of \(C\). The full tuple generates
\(A^d\) everywhere, and regularizing its maximum gives a global
smooth positive metric by \namedref{Lemma}{lem:smooth-fs}.
In the frame \(s_0\), its weight on \(\Int C\) is precisely
\(H_\varepsilon(\log|h|)\). Its curvature there is
\[
 H_\varepsilon''(\log|h|)
       \,d'\log|h|\wedge d''\log|h|
 \geq c\,d'\log|h|\wedge d''\log|h|.
\]

Now cover \(Y^\an\) by the interiors of finitely many such boxes,
using common moment presentations for the finitely many forms
\(\eta_i\) on each box. Replace the analytic moment units by
rational units as above. The symmetric coefficient matrices of
all the \(\eta_i\) are bounded on these compact boxes.
The construction for each moment unit supplies a global positive
curvature controlling its coordinate square. Suitable positive
integer multiples of these curvatures therefore dominate both
signs of every coefficient matrix on the corresponding box.
Their sum is a global smooth positive curvature, dominates every
\(\eta_i\) on the finite cover, and is the curvature of a positive
tensor power of \(A\). If the list of required coordinate squares
is empty, add any smooth positive approximation of a pure-FS
metric on \(A\). This gives the asserted ample line \(H\) and
metric \(\phi\).
\end{proof}

\begin{proposition}[Comparison on a fixed smooth test]
\label{prop:metric-comparison}
Let \(Y/K\) be smooth projective of dimension \(n\geq1\). For
\(1\leq i\leq n\), let \(p_i,q_i\) be globally smoothly
approximable metrics on the same specified line bundle \(L_i\).
Fix a real smooth function \(\varphi\) on \(Y^\an\).
Choose an ample \(H_\varphi\) and a global smooth positive
curvature \(\omega_\varphi\) such that
\[
 -\omega_\varphi\leq\ddc\varphi\leq\omega_\varphi,
\]
as in \namedref{Lemma}{lem:global-calibration}. Then
\begin{equation}\label{eq:metric-comparison}
 \begin{split}
 &\left|\int_{Y^\an}\varphi\left(
      \bigwedge_{i=1}^n c_1(L_i,p_i)
       -\bigwedge_{i=1}^n c_1(L_i,q_i)\right)\right|\\
 &\hspace{1em}\leq
  \sum_{i=1}^n\|p_i-q_i\|_\infty\,
      \deg_K\!\left(H_\varphi\prod_{h\ne i}L_h\right).
 \end{split}
\end{equation}
The norm is the uniform norm of the scalar difference of two
weights on the same line. The line \(H_\varphi\) is fixed once
\(\varphi\) and its calibration have been chosen; the degree
constants are independent of all the metrics \(p_i,q_i\).
The assertion also holds for rational line bundles after clearing
and dividing tensor denominators.
\end{proposition}

\begin{proof}
Put \(\alpha_i=c_1(L_i,p_i)\), \(\beta_i=c_1(L_i,q_i)\), and
\(T_i=\bigwedge_{h<i}\alpha_h\wedge\bigwedge_{h>i}\beta_h\).
Telescoping and current duality give
\[
 \int\varphi\left(\bigwedge_i\alpha_i-\bigwedge_i\beta_i\right)
 =\sum_i\int(p_i-q_i)\,\ddc\varphi\wedge T_i.
\]
Each \(T_i\) is positive and closed, so
\(\lvert\ddc\varphi\wedge T_i\rvert\leq\omega_\varphi\wedge T_i\).
By \namedref{Lemma}{lem:degree}, the latter measure has mass
\(\deg_K(H_\varphi\prod_{h\ne i}L_h)\). Estimating each summand by
its uniform coefficient norm proves \eqref{eq:metric-comparison}.
Multilinearity gives the rational-line version.
\end{proof}

\section*{Notation}
\renewcommand{\arraystretch}{1.06}
\begin{longtable}{@{}>{\raggedright\arraybackslash}p{0.18\textwidth}>{\raggedright\arraybackslash}p{0.50\textwidth}>{\raggedright\arraybackslash}p{0.24\textwidth}@{}}
\toprule
Symbol & Meaning & Reference\\
\midrule
\endhead
\(K,R,\pi,\lambda\)& Fixed completion, valuation ring, uniformizer,
and \(-\log|\pi|\)&\namedref{Section}{sec:introduction}\\
\(\ddc,\Ac^{p,q}\)& Original CLD differentials and compact smooth
test forms&\namedref{Section}{sec:introduction}\\
\(\Ploc,\Dloc\)& Local smooth positive approximation and its
difference class&\namedref{Definition}{def:approachable}\\
\(M_u(T)\)& The operation \(\ddc(uT)\) for approachable \(u\)&
\namedref{Proposition}{prop:fixed-potential}\\
\(\mathcal B_S\)& Mixed product with a fixed positive closed
background&\eqref{eq:recursive-product}\\
\(Q_C,\Gamma_C\)& Curve Dirichlet energy and its energy measure&
\namedref{Lemma}{lem:curve-green}\\
\(T_a\)& Normalized smooth complete-intersection curve currents&
\namedref{Proposition}{prop:curve-approximation}\\
\(p,r,g\)& Lower approachable metric, upper curve-positive metric,
and \(g=r-p\)&\namedref{Theorem}{thm:contact}\\
\(P_L(r)\)& Semipositive envelope of an obstacle metric&
\namedref{Section}{sec:envelopes}\\
\(V_l,\mathfrak j_m\)& Admissible sections and sums of divisor test ideals&
\namedref{Proposition}{prop:test-ideal-estimate}\\
\(\cE,\Delta_{\cW}\)& Vertical obstacle ideal and canonical
discrepancy boundary&\namedref{Proposition}{prop:test-ideal-estimate}\\
\(\omega_\varphi,H_\varphi\)& Fixed positive calibration and its ample line&
\namedref{Proposition}{prop:metric-comparison}\\
\(b_v,\overline\Delta_\partial\)& Boundary Green function and full
boundary datum&\namedref{Definition}{def:boundary}\\
\(\mu_v^{\mathrm{YZ}}\)& Prescribed signed adelic mixed measure&
\eqref{eq:adelic-expansion}\\
\(R_{h,\varepsilon}\)& Convex regularized maximum&
\namedref{Proposition}{prop:regularized-max}\\
\bottomrule
\end{longtable}

\section*{Declaration on the use of generative AI}
Generative AI was used in developing and drafting mathematical
arguments and in revising the exposition of this manuscript.

\begin{thebibliography}{99}
\bibitem{Ber93}
V.~G.~Berkovich,
\emph{\'Etale cohomology for non-Archimedean analytic spaces},
Publications Math\'ematiques de l'IH\'ES \textbf{78} (1993), 5--161.
\bibitem{Stacks}
The Stacks Project Authors,
\emph{The Stacks Project},
\href{https://stacks.math.columbia.edu/tag/0FUE}{Tag 0FUE},
Lemma 37.17.1.
\bibitem{CLD12}
A.~Chambert-Loir and A.~Ducros,
\emph{Formes diff\'erentielles r\'eelles et courants sur les espaces
de Berkovich}, \href{https://arxiv.org/abs/1204.6277v1}{arXiv:1204.6277v1},
2012.
\bibitem{GJRforms}
W.~Gubler, P.~Jell, and J.~Rabinoff,
\emph{Forms on Berkovich spaces based on harmonic tropicalizations},
\href{https://arxiv.org/abs/2111.05741v3}{arXiv:2111.05741v3}, 2025.
\bibitem{GJRcurves}
W.~Gubler, P.~Jell, and J.~Rabinoff,
\emph{Dolbeault cohomology of graphs and Berkovich curves},
\href{https://arxiv.org/abs/2111.05747v1}{arXiv:2111.05747v1}, 2021.
\bibitem{BMP}
B.~Bhatt, L.~Ma, Z.~Patakfalvi, K.~Schwede, K.~Tucker, J.~Waldron,
and J.~Witaszek, with an appendix by R.~Datta,
\emph{Test ideals in mixed characteristic: a unified theory up to
perturbation},
\href{https://arxiv.org/abs/2401.00615v2}{arXiv:2401.00615v2}.
\bibitem{BE}
S.~Boucksom and D.~Eriksson,
\emph{Spaces of norms, determinant of cohomology and Fekete points
in non-Archimedean geometry},
\href{https://arxiv.org/abs/1805.01016v3}{arXiv:1805.01016v3}, 2021.
\bibitem{BGM}
S.~Boucksom, W.~Gubler, and F.~Martin,
\emph{Non-Archimedean volumes of metrized nef line bundles},
\'Epijournal de G\'eom\'etrie Alg\'ebrique \textbf{5} (2021),
Article 13;
\href{https://arxiv.org/abs/2011.06986v3}{arXiv:2011.06986v3}.
\bibitem{YZ}
X.~Yuan and S.-W.~Zhang,
\emph{Adelic line bundles on quasi-projective varieties},
\href{https://arxiv.org/abs/2105.13587v9}{arXiv:2105.13587v9}.
\bibitem{CLD25}
A.~Chambert-Loir and A.~Ducros,
\emph{Formes diff\'erentielles r\'eelles et courants sur les espaces
de Berkovich},
\href{https://arxiv.org/abs/1204.6277v2}{arXiv:1204.6277v2}, 2025.
\bibitem{Wanner}
V.~Wanner,
\emph{Comparison of two notions of subharmonicity on non-archimedean
curves},
\href{https://arxiv.org/abs/1801.04713v3}{arXiv:1801.04713v3}, 2018.
\bibitem{BGJK}
J.~I.~Burgos Gil, W.~Gubler, P.~Jell, and K.~K\"unnemann,
\emph{Pluripotential theory for tropical toric varieties and
non-archimedean Monge--Amp\`ere equations},
\href{https://arxiv.org/abs/2102.07392v1}{arXiv:2102.07392v1}, 2021.
\bibitem{FGK}
Y.~Fang, W.~Gubler, and K.~K\"unnemann,
\emph{On the non-archimedean Monge--Amp\`ere equation in mixed
characteristic},
\href{https://arxiv.org/abs/2203.12282v3}{arXiv:2203.12282v3}, 2025.
\bibitem{GRpl}
W.~Gubler and J.~Rabinoff,
\emph{On semipositive piecewise linear functions in non-archimedean
analytic geometry},
\href{https://arxiv.org/abs/2503.05513v1}{arXiv:2503.05513v1}, 2025.
\bibitem{GRpsh}
W.~Gubler and J.~Rabinoff,
\emph{Classically psh and pluriharmonic functions on Berkovich spaces},
\href{https://arxiv.org/abs/2506.13548v2}{arXiv:2506.13548v2}, 2025.
\bibitem{Wu}
L.~Wu,
\emph{On the local theory of non-archimedean psh functions},
\href{https://arxiv.org/abs/2609.27610v1}{arXiv:2609.27610v1}, 2026.
\end{thebibliography}
\end{document}
